\chapter{RESULTS AND DISCUSSION}\label{chapter-v.-results-and-discussion}

% Chapter opens at §5.0: \chapter reset the section counter to 0; nudging it to -1 makes the first
% \section (Statistical Methodology) number as 5.0, matching the prose scheme (ladder = §5.0.1,
% EXP-005 = §5.1, ..., EXP-010 = §5.10), with the re-added §5.11/§5.12/§5.13 following. Localized:
% ch6's \chapter resets the counter, so subsequent chapters are unaffected.
\setcounter{section}{-1}

\section{Statistical Methodology and Reproducibility}\label{statistical-methodology-and-reproducibility}

All headline statistics in §§5.1--5.4, §5.7, §5.9 are computed across \textbf{20 independent seeds (1--20)} with the actuarial simulator run for \textbf{N = 5,000 rounds} per seed. Two supplementary experiments use \textbf{10 seeds per cell} by design: §5.8 (sensitivity sweep across α, adverse-selection factor, and elasticity slope --- 11 cells × 10 seeds) and §5.10 (cold-start analysis at four horizons, each requiring a fresh XGBoost refit). The divergence is documented in the respective source files (\texttt{exp\_012\_sensitivity\_analysis.py:48}, \texttt{exp\_010\_cold\_start\_analysis.py:253}) and motivated by computational cost; results from the 10-seed cells are reported with the same bootstrap/Wilcoxon machinery and noted as such in each section. Tables report \textbf{mean ± standard deviation} with \textbf{95 \% bootstrap confidence intervals} in brackets. Pairwise comparisons between algorithms use the \textbf{paired Wilcoxon signed-rank test} with \textbf{Bonferroni correction} for multiple comparisons; effect sizes are reported as \textbf{Cohen's d}. The primary illustrative seed (used for trajectory figures) is SEED = 42; figures are produced from this seed unless noted.

A reviewer can reproduce any reported number by running \texttt{python\ healthrl/experiments/exp\_005\_underwriting\_convergence.py} (or the corresponding \texttt{exp\_006}, \texttt{exp\_007}, \texttt{exp\_008}) on Python 3.11 with the pinned dependencies in \texttt{requirements.txt}. Each script exits with code 0 if the pre-registered pass criteria are satisfied.

A central methodological commitment of this chapter is \textbf{baseline completeness}. Earlier drafts benchmarked the learners against a single hand-tuned XGBoost rule. Following the principle that a proposed method must be measured against the \emph{best simple alternative} --- not merely a convenient weak one --- the evaluation reported here adds a complete \textbf{baseline ladder}: a uniform-random policy, the three constant policies (AlwaysSTANDARD, AlwaysRATED, AlwaysDECLINE), the frozen Static-XGB rule, and a fully-supervised linear classifier fit on oracle-optimal labels (LogisticOracle, an \emph{in-sample optimistic ceiling}, not a deployable policy). All policies are evaluated under \textbf{identical common random numbers} within each seed, so every comparison sees the same applicants, acceptance draws, and claims noise. The pre-registered pass/fail criteria are reported honestly, \emph{including those that were falsified} (§5.3, §5.7) --- a falsified pre-registration is a result, not a failure to hide.

\subsection{The Complete Baseline Ladder and the Headline Question}\label{the-complete-baseline-ladder-and-the-headline-question}

\subsubsection{The headline question}\label{the-headline-question}

The operative question of this thesis is not \emph{``does a contextual bandit beat a hand-tuned rule?''} --- a low bar a learner clears almost by construction --- but \textbf{``does a contextual bandit beat the best \emph{simple} policy available on this problem?''} Answering it requires a complete baseline ladder. Table 9 reports that ladder under the as-published (``ORIGINAL'') reward model, 20 seeds, N = 5,000 rounds, common random numbers. (This promotes the 5-seed first pass in EXP-014 to the chapter's standard 20-seed protocol; the reimplemented reward harness reproduces the headline LinUCB reward of EXP-007's original run to within rounding --- \$91,864 on the harness vs \$91,947 from EXP-007's own pipeline --- confirming the two pipelines agree. Table 12 (§5.3) reports the harness values for LinUCB, LinTS, and Static XGB, so that every ladder row shares the same common random numbers. EXP-005's own harness reports a slightly lower \$90,540 for LinUCB (§5.1); the \textasciitilde1.4 \% difference is between EXP-005's pipeline and the regression-exact reward-sensitivity harness used for the ladder, and is immaterial to every comparison in this chapter.)

\captionof{table}{Table 9 --- The baseline ladder (ORIGINAL reward, N = 5,000, 20 seeds, mean cumulative reward {[}95 \% bootstrap CI{]})}\label{tab:5.0.1}

{\def\LTcaptype{none} % do not increment counter
\begin{longtable}[]{@{}
  >{\raggedright\arraybackslash}p{(\linewidth - 6\tabcolsep) * \real{0.1739}}
  >{\raggedleft\arraybackslash}p{(\linewidth - 6\tabcolsep) * \real{0.4130}}
  >{\raggedleft\arraybackslash}p{(\linewidth - 6\tabcolsep) * \real{0.2826}}
  >{\raggedright\arraybackslash}p{(\linewidth - 6\tabcolsep) * \real{0.1304}}@{}}
\toprule\noalign{}
\begin{minipage}[b]{\linewidth}\raggedright
Policy
\end{minipage} & \begin{minipage}[b]{\linewidth}\raggedleft
Cumulative reward
\end{minipage} & \begin{minipage}[b]{\linewidth}\raggedleft
\% of Oracle
\end{minipage} & \begin{minipage}[b]{\linewidth}\raggedright
Note
\end{minipage} \\
\midrule\noalign{}
\endhead
\bottomrule\noalign{}
\endlastfoot
AlwaysDECLINE & -50,000 (deterministic) & -39.6 \% & constant decline floor \\
Random & 1,980 {[}739, 3,209{]} & 1.6 \% & uniform action floor \\
AlwaysSTANDARD & 34,684 {[}33,694, 35,871{]} & 27.5 \% & conservative constant \\
Static XGB (hand-tuned rule) & 72,206 {[}70,125, 74,352{]} & 57.1 \% & the original sole baseline \\
\textbf{LinUCB} (proposed) & \textbf{91,864 {[}89,139, 94,702{]}} & \textbf{72.7 \%} & rank-2 of the admissible policies \\
\textbf{LinTS} (proposed) & \textbf{93,723 {[}91,105, 96,481{]}} & \textbf{74.2 \%} & rank-1 of the admissible policies \\
AlwaysRATED & 122,287 {[}119,528, 124,996{]} & 96.8 \% & best constant; \emph{inadmissible} (§5.3, §5.11) \\
LogisticOracle ‡ & 123,977 {[}121,794, 126,126{]} & 98.1 \% & full-supervision ceiling, \emph{not deployable} \\
Oracle (upper bound) & 126,351 {[}124,187, 128,428{]} & 100 \% & --- \\
\end{longtable}
}

‡ LogisticOracle is a logistic classifier fit \textbf{in-sample on the oracle-optimal labels}. It therefore requires the very answer it is predicting and is \textbf{not a deployable baseline} --- it is an \emph{optimistic static ceiling}. Its role is diagnostic (see below).

Two terms are used in a fixed sense throughout this chapter. A policy is \textbf{deployable} if it can be implemented in production from observable applicant features alone --- AlwaysRATED qualifies; Oracle and LogisticOracle, which require knowledge of the optimal action, do not. A policy is \textbf{admissible} if it is deployable \emph{and} commercially and regulatorily viable as an underwriting rule --- AlwaysRATED, which rates 100 \% of applicants at a flat +25 \% loading, is deployable but \emph{inadmissible} (§6.2).

\fg{0.85\textwidth}{Figure 6. The baseline ladder (mean ± 95 \% bootstrap CI). LinUCB and LinTS lead every admissible policy, outperforming Static XGB by +27.2 \% (LinUCB) and +29.8 \% (LinTS) on the ladder-harness basis (Table~9)}{images/fig_exp_014_baseline_ladder.pdf}

\subsubsection{The verdict, stated up front}\label{the-verdict-stated-up-front}

\textbf{The contextual bandit beats every admissible alternative on this problem.} LinUCB and LinTS lead the admissible-policy rankings in Table 9, outperforming the frozen Static-XGB rule by \textbf{+27.2 \%} and \textbf{+29.8 \%} respectively (LinUCB: p \textless{} 0.001, d = 2.98, large effect; §5.1, §5.3). Both bandits also clear every other admissible baseline --- Epsilon-Greedy, AlwaysSTANDARD, and Random --- by large, statistically significant margins. \emph{Online adaptive estimation beats every admissible frozen alternative.}

A post-hoc robustness analysis identified an inadmissible constant policy (AlwaysRATED --- rating 100 \% of applicants at a +25 \% loading, which is neither commercially nor regulatorily feasible) as a theoretical performance ceiling above the bandits (Table 9). This finding is quantified and interpreted in §6.2.

\section{EXP-005: Convergence Validation}\label{exp-005-convergence-validation}

\fg{0.85\textwidth}{Figure 7. Cumulative reward curves for LinUCB and Static XGB baseline (mean of 20 seeds, shaded band = 95\% bootstrap confidence interval)}{images/fig_reward_curves.pdf}

\subsection{Cumulative Reward and Regret}\label{cumulative-reward-and-regret}

Table 10 reports the cumulative reward and average regret for LinUCB against the Oracle ceiling and the Static XGBoost baseline over N = 5,000 rounds, averaged over 20 independent seeds.

\captionof{table}{Table 10 --- EXP-005 cumulative performance (N = 5,000 rounds, 20 seeds, mean ± std {[}95\% bootstrap CI{]})}\label{tab:5.1.1}

{\def\LTcaptype{none} % do not increment counter
\begin{longtable}[]{@{}
  >{\raggedright\arraybackslash}p{(\linewidth - 6\tabcolsep) * \real{0.1778}}
  >{\centering\arraybackslash}p{(\linewidth - 6\tabcolsep) * \real{0.1778}}
  >{\centering\arraybackslash}p{(\linewidth - 6\tabcolsep) * \real{0.1778}}
  >{\centering\arraybackslash}p{(\linewidth - 6\tabcolsep) * \real{0.4667}}@{}}
\toprule\noalign{}
\begin{minipage}[b]{\linewidth}\raggedright
Metric
\end{minipage} & \begin{minipage}[b]{\linewidth}\centering
Oracle
\end{minipage} & \begin{minipage}[b]{\linewidth}\centering
LinUCB
\end{minipage} & \begin{minipage}[b]{\linewidth}\centering
Static XGB Baseline
\end{minipage} \\
\midrule\noalign{}
\endhead
\bottomrule\noalign{}
\endlastfoot
Cumulative Reward & 126,804 ± 4,724 {[}124,857, 128,846{]} & \textbf{90,540 ± 5,382 {[}88,287, 92,886{]}} & 72,292 ± 4,120 {[}70,646, 74,136{]} \\
Avg Regret (last 500 rounds) & -3.06 ± 2.51 {[}-4.04, -1.88{]} & \textbf{2.20 ± 2.77 {[}1.00, 3.37{]}} & 9.01 ± 3.87 {[}7.30, 10.62{]} \\
\end{longtable}
}

LinUCB accumulates a mean of \textbf{90,540} in total reward over 5,000 rounds, exceeding the Static XGB baseline by \textbf{+18,248 (+25.2 \%)} and recovering roughly 72 \% of the Oracle ceiling (126,804). This confirms that \emph{online adaptive estimation beats a frozen rule} --- the narrow claim that survives the baseline-ladder analysis (§5.0.1). It does \textbf{not}, however, establish that LinUCB learns a strong policy in absolute terms: the same 72 \%-of-oracle reward leaves it a full tier below the trivial AlwaysRATED constant (96.8 \% of oracle; §5.0.1, §5.3). The late-round average regret is more diagnostic of the gap to the \emph{rule}: LinUCB pays only 2.20 per round in the final 500 rounds --- within 5 of the Oracle floor --- while the Static XGB baseline continues to misallocate underwriting actions at four times that rate (9.01/round). The gap to Static XGB widens monotonically across the run, but §5.13 shows that what the bandit converges \emph{to} is a steady-state policy whose per-round reward remains below the constant's --- so the convergence demonstrated here is convergence to a \emph{profitable-but-suboptimal} policy, not to the best available one.

\textbf{Statistical inference (paired, 20 seeds).} The reward difference is highly significant under a paired Wilcoxon signed-rank test (p \textless{} 0.001, mean diff +18,248, 95\% CI {[}+15,595, +20,858{]}, Cohen's d = 2.98, large effect). The regret difference is equally robust (p \textless{} 0.001, mean diff -6.81, 95\% CI {[}-8.60, -4.90{]}, Cohen's d = -1.59, large effect). Both pairwise comparisons remain significant after Bonferroni correction for the two tests.

\fg{0.85\textwidth}{Figure 8. Evolution of action distribution over 5,000 rounds. Entropy decreases from 1.31 (uniform-like exploration) to 1.10 (exploitation regime)}{images/fig_action_evolution.pdf}

\subsection{Action Distribution Evolution and Entropy}\label{action-distribution-evolution-and-entropy}

Action entropy provides a seed-invariant measure of exploration intensity. Early in the run (rounds 1--500), LinUCB distributes its four actions relatively uniformly: mean entropy is 1.314 ± 0.017 nats, close to but below the uniform-distribution maximum of \(\log 4 ≈ 1.386\) nats. By rounds 4,500--5,000, mean entropy has fallen to 1.105 ± 0.036 nats --- a statistically significant decrease that reflects the bandit's progression from exploration to exploitation. Across all 20 seeds the late entropy is strictly less than the early entropy, satisfying the convergence pass criterion.

Qualitatively, this entropy decrease is driven principally by the REFER arm: in the early phase REFER is over-selected because its prior is uninformative and its UCB bonus is correspondingly inflated; once the bandit has observed enough (context, action, reward) triples to estimate reliable feature--reward coefficients for the dominant applicant archetypes, REFER selection collapses and the policy concentrates on STANDARD, RATED, and DECLINE in proportions consistent with the underlying actuarial risk distribution.

\subsection{Learning Curve Analysis}\label{learning-curve-analysis}

The reward curve separates from the Static XGB baseline within the first few hundred rounds, after which LinUCB's cumulative reward grows at a consistently steeper slope until the gap stabilises at roughly 18,000 by the end of the run. The corresponding regret curve grows approximately as \(O(\sqrt{t})\) in the early phase --- consistent with the theoretical LinUCB bound of \(\tilde{O}(d\sqrt{T})\) for a \(d\)-dimensional feature space \citep{Li2010} --- before flattening as the bandit's posterior tightens. §5.6 reports a direct empirical validation of this T-dependence: a log-log fit of mean cumulative regret across 20 seeds recovers a slope of 0.564 ± 0.142 (mean-curve slope 0.572, R² = 0.99), bracketing the theoretical 0.5.

Action accuracy against the Oracle is \textbf{37.8\% ± 3\%} in the last 500 rounds. This is well above the 30\% pass threshold but far from 100\%, and the interpretation deserves emphasis: the bandit is not converging to the Oracle's exact policy. Instead it is discovering a \emph{different but profitable} policy that achieves 72\% of the Oracle's reward. This pattern reflects a limited-feature regime: because the linear reward model in \texttt{expected\_rewards()} uses the same 34-feature context that the bandit observes --- while the Oracle has access to the per-sample stochastic noise realisation --- the bandit cannot match the Oracle's action choices case-by-case. Instead it learns an alternative policy that is near-optimal \emph{within the available feature subspace}. From an actuarial perspective the bandit learns an alternative policy that is \emph{locally} sensible given observable features; but ``maximises observable reward'' overstates it --- a constant policy and a fully-supervised linear model on the same features both score higher (§5.0.1, §5.11). The 37.8 \% Oracle-agreement should be read as \emph{the bandit settling into a distinct, lower-reward policy}, the structural character of which is examined in §5.13.

\subsection{Pass/Fail Assessment}\label{passfail-assessment}

Five pre-registered criteria were evaluated on 20 seeds:

\begin{enumerate}
\def\labelenumi{\arabic{enumi}.}
\tightlist
\item
  \textbf{LinUCB cumulative reward significantly \textgreater{} Static XGB} --- +18,248, p \textless{} 0.001, Wilcoxon. \textbf{PASSED.}
\item
  \textbf{LinUCB avg regret (last 500) significantly \textless{} Static XGB} --- -6.81, p \textless{} 0.001, Wilcoxon. \textbf{PASSED.}
\item
  \textbf{Action entropy decreases (exploration → exploitation)} --- 1.314 → 1.105. \textbf{PASSED.}
\item
  \textbf{Action accuracy (last 500) \textgreater{} 30\%} --- 37.8\% (different but profitable policy). \textbf{PASSED.}
\item
  \textbf{Oracle action accuracy = 100\%} (sanity check). \textbf{PASSED.}
\end{enumerate}

\section{EXP-006: Fairness Audit}\label{exp-006-fairness-audit}

\fg{0.85\textwidth}{Figure 9. Converged-phase regional approval rates (20-seed means with 95\% CI) and the EEOC four-fifths (80\% of maximum) overlay. All eight regions clear the threshold; per-seed parity ratios are reported in Table 11.}{images/fig_fairness_region.pdf}

\subsection{Sliding-Window PSI}\label{sliding-window-psi}

Rather than computing a single PSI between the full applicant pool and the final approved portfolio, EXP-006 evaluates PSI on a 500-round sliding window against a fixed reference window (the first 500 rounds of each seed). This captures whether the bandit's learning trajectory introduces \emph{temporal} demographic concentration, not only end-state concentration. The maximum sliding-window PSI across the run is the relevant guardrail metric for live deployment.

\begin{itemize}
\tightlist
\item
  \textbf{Region} --- maximum sliding-window PSI = 0.0821 ± 0.0213 (95\% CI {[}0.0727, 0.0913{]}); final-window PSI = 0.0424 ± 0.0233. Both values are in the \textbf{GREEN} zone (\textless{} 0.10), confirming that LinUCB never induces a regionally concentrated approval pattern at any point during the run.
\item
  \textbf{Occupation} --- maximum sliding-window PSI = 0.1225 ± 0.0793 (95\% CI {[}0.0953, 0.1613{]}); final-window PSI = 0.0710 ± 0.0299. The maximum value sits at the lower edge of the \textbf{AMBER} band on average but well below the \textbf{RED} threshold (\textless{} 0.25); the final-window value is comfortably \textbf{GREEN}. This indicates a transient occupational concentration during exploration that resolves itself as the bandit converges.
\end{itemize}

Figure 10 traces both sliding-window PSI series across the full run against the GREEN/AMBER/RED guardrail bands, showing the occupation series' transient rise toward the AMBER zone during early exploration and its return into GREEN as the policy converges.

\fg{0.85\textwidth}{Figure 10. Sliding-window Population Stability Index for region and occupation across the 5,000-round run (20-seed mean; shaded band = 95\% bootstrap CI) against the GREEN ($<0.10$), AMBER ($0.10$--$0.25$), and RED ($>0.25$) guardrail thresholds. Region PSI stays within GREEN throughout; the occupation series rises toward the GREEN/AMBER boundary during early exploration---its 95\% confidence band extends into the lower AMBER zone---then eases back as the policy converges, never approaching RED.}{images/fig_psi_timeseries.pdf}

\subsection{Approval-Rate Parity (EEOC 4/5 Rule)}\label{approval-rate-parity-eeoc-45-rule}

Approval rates (proportion of STANDARD + RATED decisions) were computed on the converged phase (rounds 3,000--4,999) for each demographic group. The pre-registered constraint is the U.S. EEOC ``four-fifths'' rule: the minimum group approval rate should be at least 80\% of the maximum group's rate.

\captionof{table}{Table 11 --- EXP-006 converged-phase approval rates (rounds 3,000--4,999, 20 seeds)}\label{tab:5.2.1}

{\def\LTcaptype{none} % do not increment counter
\begin{longtable}[]{@{}
  >{\raggedright\arraybackslash}p{(\linewidth - 8\tabcolsep) * \real{0.1279}}
  >{\centering\arraybackslash}p{(\linewidth - 8\tabcolsep) * \real{0.1628}}
  >{\centering\arraybackslash}p{(\linewidth - 8\tabcolsep) * \real{0.1628}}
  >{\centering\arraybackslash}p{(\linewidth - 8\tabcolsep) * \real{0.2791}}
  >{\centering\arraybackslash}p{(\linewidth - 8\tabcolsep) * \real{0.2674}}@{}}
\toprule\noalign{}
\begin{minipage}[b]{\linewidth}\raggedright
Dimension
\end{minipage} & \begin{minipage}[b]{\linewidth}\centering
Min Approval
\end{minipage} & \begin{minipage}[b]{\linewidth}\centering
Max Approval
\end{minipage} & \begin{minipage}[b]{\linewidth}\centering
Parity Ratio (Min/Max)
\end{minipage} & \begin{minipage}[b]{\linewidth}\centering
EEOC 4/5 Rule (≥ 80\%)
\end{minipage} \\
\midrule\noalign{}
\endhead
\bottomrule\noalign{}
\endlastfoot
Region & 0.6622 ± 0.0493 & 0.7725 ± 0.0315 & \textbf{85.72\%} & \textbf{PASS} \\
Occupation & 0.6803 ± 0.0303 & 0.7549 ± 0.0301 & \textbf{90.12\%} & \textbf{PASS} \\
\end{longtable}
}

Both dimensions satisfy the four-fifths rule with comfortable margin. The rural--urban approval gap and the cross-occupational gap reflect actuarially justified differences in the underlying mortality and morbidity distribution of the Cambodia dataset --- rural and high-labour-intensity occupations carry higher prevalence of Hepatitis B, TB co-morbidities, and physical-injury risk --- rather than proxy discrimination on region or occupation as demographic characteristics per se.

\fg{0.85\textwidth}{Figure 11. Converged-phase occupational approval rates (20-seed means with 95\% CI) and the EEOC four-fifths (80\% of maximum) overlay. All seven occupations clear the threshold; per-seed parity ratios are reported in Table 11.}{images/fig_fairness_occupation.pdf}

\subsection{Permutation Independence Tests}\label{permutation-independence-tests}

To assess whether action allocation depends statistically on demographic group, action labels were permuted within each seed and the empirical distribution of a test statistic compared to the permuted null. Bonferroni-corrected significance threshold is \(\alpha / 2 = 0.025\) for the two tests.

\begin{itemize}
\tightlist
\item
  \textbf{Region} --- p = 0.1322 ± 0.1509 (95\% CI {[}0.0733, 0.2017{]}). The null of action--region independence cannot be rejected; there is no significant regional bias in the learned policy.
\item
  \textbf{Occupation} --- p \textless{} 0.001 (95\% CI {[}0.0000, 0.0000{]}). The null of action--occupation independence is rejected, indicating a statistically significant association between action and occupation.
\end{itemize}

The occupational association is statistically significant but practically small: the parity ratio is 90.12\% (well above the 80\% EEOC threshold) and the maximum sliding-window PSI is 0.1225 (within the GREEN/AMBER boundary). This pattern is expected: different occupational groups have genuinely different actuarial risk profiles, so a policy maximising expected reward will allocate actions in occupation-dependent proportions. The fairness guardrail is not ``no statistical dependence'' but ``no disparate-impact magnitude beyond the regulatory threshold,'' and that guardrail is satisfied.

\subsection{Pass/Fail Assessment}\label{passfail-assessment-1}

\begin{enumerate}
\def\labelenumi{\arabic{enumi}.}
\tightlist
\item
  \textbf{Region: min/max approval ratio ≥ 80\% (EEOC 4/5 rule)} --- 85.72\%. \textbf{PASSED.}
\item
  \textbf{Occupation: min/max approval ratio ≥ 80\% (EEOC 4/5 rule)} --- 90.12\%. \textbf{PASSED.}
\item
  \textbf{Region: max sliding-window PSI \textless{} 0.25} --- 0.0821. \textbf{PASSED.}
\item
  \textbf{Occupation: max sliding-window PSI \textless{} 0.25} --- 0.1225. \textbf{PASSED.}
\item
  \textbf{Region: permutation test p \textgreater{} 0.025 (action--region independence)} --- p = 0.1322 (20-seed mean, 95\% CI {[}0.0733, 0.2017{]}); fails to reject the null. \textbf{PASSED.}
\item
  \textbf{Occupation: permutation test p \textgreater{} 0.025 (action--occupation independence)} --- p \textless{} 0.001 (95\% CI {[}0.000, 0.000{]}); the null of action--occupation independence is rejected. \textbf{FAILED.} The association is statistically significant but practically small: the parity ratio is 90.12\% (above the 80\% EEOC threshold) and the maximum PSI is 0.1225 (within the GREEN/AMBER boundary), satisfying the regulatory guardrail. The dependence reflects actuarially justified risk differences across occupational groups; the fairness criterion is ``no disparate impact beyond the regulatory threshold,'' not ``no statistical dependence,'' and that criterion is satisfied. This criterion is restored from the pre-registered list and marked FAILED in the interest of honesty; the practical finding is that the PSI/EEOC apparatus is operating as a monitor, not an enforcer, exactly as specified in §4.8.
\end{enumerate}

A scope note for honesty: the PSI/EEOC apparatus demonstrated here is a \textbf{monitor}, not an \textbf{enforcer} --- it detects demographic concentration but does not constrain the policy to avoid it. An enforcement mechanism (e.g., a constrained-action layer) is identified as future work in Chapter VI. This does not weaken the present finding (no threshold is breached), but it bounds the claim to \emph{measurement}. The fairness result is otherwise \textbf{orthogonal to the reward-design question and survives intact}: it is one of the chapter's genuine, unaffected contributions.

\section{EXP-007: Benchmark Comparison}\label{exp-007-benchmark-comparison}

\fg{0.85\textwidth}{Figure 12. Cumulative regret curves across all four algorithms (mean of 20 seeds, shaded band = 95\% CI)}{images/fig_regret_curves.pdf}

\subsection{Final Rankings}\label{final-rankings}

Table 12 ranks all four algorithms by mean cumulative reward and mean cumulative regret over 5,000 rounds, averaged across 20 independent seeds.

\captionof{table}{Table 12 --- EXP-007 benchmark over the complete baseline ladder (ORIGINAL reward, N = 5,000, 20 seeds, mean cumulative reward {[}95 \% CI{]})}\label{tab:5.3.1}

Reward is the primary metric because the realized-regret sign for the constant policies is distorted by claims noise (§5.11). LinUCB, LinTS, and Static XGB are reported on the same 20-seed common-random-number harness as the constant policies and LogisticOracle (§5.11), which reproduces EXP-007's own-run values to within rounding (LinUCB \$91,864 vs \$91,947; LinTS \$93,723 vs \$93,572; Oracle 126,351 vs 126,318); Epsilon-Greedy, which is not part of the harness, uses EXP-007's own run.

{\def\LTcaptype{none} % do not increment counter
\begin{longtable}[]{@{}
  >{\centering\arraybackslash}p{(\linewidth - 6\tabcolsep) * \real{0.0909}}
  >{\raggedright\arraybackslash}p{(\linewidth - 6\tabcolsep) * \real{0.2000}}
  >{\raggedleft\arraybackslash}p{(\linewidth - 6\tabcolsep) * \real{0.3455}}
  >{\raggedleft\arraybackslash}p{(\linewidth - 6\tabcolsep) * \real{0.3636}}@{}}
\toprule\noalign{}
\begin{minipage}[b]{\linewidth}\centering
Rank (deployable)
\end{minipage} & \begin{minipage}[b]{\linewidth}\raggedright
Policy
\end{minipage} & \begin{minipage}[b]{\linewidth}\raggedleft
Cumulative reward
\end{minipage} & \begin{minipage}[b]{\linewidth}\raggedleft
\% of best constant
\end{minipage} \\
\midrule\noalign{}
\endhead
\bottomrule\noalign{}
\endlastfoot
--- & Oracle (ceiling) & 126,351 {[}124,187, 128,428{]} & --- \\
--- & LogisticOracle ‡ (full-supervision ceiling) & 123,977 {[}121,794, 126,126{]} & 101.4 \% \\
\textbf{1} & \textbf{AlwaysRATED} (trivial constant; \emph{inadmissible}, §5.0.1) & \textbf{122,287 {[}119,528, 124,996{]}} & \textbf{100 \%} \\
2 & LinTS & 93,723 {[}91,105, 96,481{]} & 76.6 \% \\
3 & LinUCB & 91,864 {[}89,139, 94,702{]} & 75.1 \% \\
4 & Epsilon-Greedy & 76,441 {[}73,767, 79,107{]} & 62.5 \% \\
5 & Static XGB & 72,206 {[}70,125, 74,352{]} & 59.1 \% \\
6 & AlwaysSTANDARD & 34,684 {[}33,694, 35,871{]} & 28.4 \% \\
7 & Random & 1,980 {[}739, 3,209{]} & 1.6 \% \\
\end{longtable}
}

‡ in-sample optimistic ceiling, not deployable (§5.0.1).

LinTS narrowly outperforms LinUCB on mean reward and regret, but the two are \textbf{not statistically distinguishable} under a paired Wilcoxon test on regret (p = 0.87, mean diff +1,626, 95\% CI {[}-1,135, +4,390{]}, Cohen's d = 0.26, small effect). The remaining pairwise comparisons are highly significant after Bonferroni correction (α / 6 = 0.0083 for the six regret tests):

\begin{itemize}
\tightlist
\item
  LinUCB vs Epsilon-Greedy (regret): -15,506, p \textless{} 0.001, d = -3.41.
\item
  LinUCB vs Static XGB (regret): -19,774, p \textless{} 0.001, d = -3.41.
\item
  LinTS vs Epsilon-Greedy (regret): -17,132, p \textless{} 0.001, d = -2.75.
\item
  LinTS vs Static XGB (regret): -21,400, p \textless{} 0.001, d = -3.89.
\item
  Epsilon-Greedy vs Static XGB (regret): -4,268, p \textless{} 0.01, d = -0.82.
\end{itemize}

The expanded ranking changes the headline. \textbf{Two findings coexist.} First, among the \emph{adaptive vs non-adaptive} comparison the original chapter drew, the result holds: LinTS and LinUCB beat the frozen Static-XGB rule and ε-Greedy by large, significant margins (regret differences p \textless{} 0.001, \textbar d\textbar{} = 2.8--3.9 after Bonferroni; the within-learner detail is unchanged and reported in §5.3.2). Second --- and decisively for the thesis --- \textbf{the contextual bandits are rank-3 among deployable policies, beaten by the trivial AlwaysRATED constant by +30--33 \%} (p \textless{} 0.0001, d = −4.5 / −5.1; §5.0.1, §5.11). The corrected takeaway is therefore: \emph{online adaptation beats a frozen rule, but neither bandit beats the best simple policy.} The choice between UCB and Thompson Sampling remains a matter of operational preference --- LinTS edges LinUCB but the two are statistically indistinguishable (paired Wilcoxon p = 0.87) --- a finding that is preserved but is now a comparison \textbf{between two tier-3 policies}.


The ablation in §5.7 sharpens this finding. A \emph{Greedy-only} variant of LinUCB (α = 0) performs the same online ridge update as standard LinUCB but suppresses the exploration bonus entirely; its mean cumulative reward is not significantly lower than full LinUCB's (paired Wilcoxon p = 0.58, Cohen's d = 0.13). The load-bearing component on the Cambodia dataset is therefore \emph{online contextual updating}, not the uncertainty-directed exploration bonus per se --- a result consistent with Bastani, Bayati \& Khosravi \citeyearpar{Bastani2021}, who prove that covariate diversity can make exploration-free greedy near-optimal.

Epsilon-Greedy (rank 3) sits apart from this picture: its ε = 0.15 uniform-random arm draws are independent of the learned coefficients and therefore do not benefit from the shrinkage-of-uncertainty mechanism that links UCB, Thompson Sampling, and α = 0 greedy through their shared linear estimator (§5.3.2).

\subsection{Algorithm-Specific Behaviour}\label{algorithm-specific-behaviour}

\textbf{LinTS (rank 1).} Thompson Sampling maintains a full Gaussian posterior over reward parameters and samples from it at each round. This produces an automatic exploration--exploitation balance: when the posterior is wide, samples are spread across arms; as the posterior narrows, samples concentrate on the optimal arm. Crucially, LinTS requires no exploration parameter tuning --- unlike LinUCB's \(\alpha\) or Epsilon-Greedy's \(\varepsilon\) --- which is an operational advantage in production settings where re-tuning after data drift is expensive.

\textbf{LinUCB (rank 2).} Upper Confidence Bound performs comparably to LinTS but requires choice of \(\alpha\). At \(\alpha = 1.0\) (Li et al., 2010 default), it accumulates slightly more regret in the early rounds than LinTS before converging to a similar late-round policy. The two algorithms reach statistically indistinguishable asymptotic performance, but UCB's parameter sensitivity means that production deployment would warrant either a tuned \(\alpha\) for the specific data distribution or a switch to Thompson Sampling.

\textbf{Epsilon-Greedy (rank 3).} With \(\varepsilon = 0.15\), the algorithm wastes 15\% of decisions on uniform random exploration regardless of how much has been learned. Unlike UCB and Thompson Sampling, it does not reduce exploration as confidence grows. This produces a nearly-linear regret curve --- the hallmark of an algorithm that never fully exploits its learned knowledge --- and its mean regret is roughly 1.8 times that of LinUCB and LinTS combined.

\textbf{Static XGB (rank 4).} The pre-trained XGBoost baseline, combined with deterministic underwriting rules, cannot adapt to the reward signals it receives. Its regret is bounded below by the systematic mismatch between its fixed decision boundaries and the actuarial reward structure. Even Epsilon-Greedy --- which explores uniformly without using any context-dependent information beyond the empirical mean --- outperforms it by a statistically significant margin (-4,268, d = -0.82, large effect), underscoring that \emph{any} online adaptation beats \emph{no} online adaptation in this environment.

The ``rank 1--4'' labels in this subsection refer to the \emph{within-learner} ordering of the four algorithms EXP-007 originally ran; in the complete ladder of Table 12 the bandits are rank-3 overall. These descriptions characterize the \emph{relative} behaviour of the four learners; none of them reaches the trivial-constant ceiling (§5.0.1), which is the subject of §5.11.

\subsection{Pass/Fail Assessment}\label{passfail-assessment-2}

\begin{enumerate}
\def\labelenumi{\arabic{enumi}.}
\tightlist
\item
  \textbf{Oracle mean regret ≤ all learning algorithms} --- Oracle -11,597 \textless{} LinTS 21,149 \textless{} LinUCB 22,774 \textless{} EpsGreedy 38,281 \textless{} StaticXGB 42,548. \textbf{PASSED.}
\item
  \textbf{LinUCB regret significantly \textless{} StaticXGB and Epsilon-Greedy} --- both p \textless{} 0.001, d = -3.41. \textbf{PASSED.}
\item
  \textbf{LinTS regret significantly \textless{} StaticXGB and Epsilon-Greedy} --- p \textless{} 0.001, d = -3.89 and d = -2.75. \textbf{PASSED.}
\item
  \textbf{LinUCB and LinTS reward significantly \textgreater{} StaticXGB} --- both p \textless{} 0.001, d \textgreater{} 3.4. \textbf{PASSED.}
\item
  \textbf{Oracle reward ≥ all learning algorithms} --- confirmed. \textbf{PASSED.}
\item
  \textbf{Pre-registered expectation \texttt{bandit\ \textgreater{}\ best\ simple\ baseline} --- FALSIFIED.} AlwaysRATED (\$122,287) \textgreater{} LinTS (\$93,723) \textgreater{} LinUCB (\$91,864), p \textless{} 0.0001. Reported in the same spirit as the falsified hypotheses in EXP-011 (§5.7): the falsification \emph{is} the scientifically informative outcome and motivates the reframing of this chapter.
\end{enumerate}

\section{EXP-008: Human-in-the-Loop Underwriting}\label{exp-008-human-in-the-loop-underwriting}

\subsection{Motivation and Design}\label{motivation-and-design}

EXP-005 through EXP-007 demonstrated that LinUCB and LinTS outperform static baselines on synthetic data where the optimal reward for every action is immediately observable. In practice, however, an underwriter deploying a contextual bandit faces a fundamentally different constraint: when the model is uncertain about a borderline applicant, it can defer to a human expert --- but human review is costly, creates queue latency, and cannot scale to every case. EXP-008 asks whether replacing the REFER arm's mathematical oracle with a real (simulated) human underwriter still produces better outcomes than a mathematical baseline while remaining cost-effective.

The experiment wraps LinUCB (α = 1.0) in a human-in-the-loop (HITL) wrapper. Whenever the bandit selects REFER, the case enters a review queue and a simulated human underwriter resolves it to a final action (STANDARD, RATED, or DECLINE). The simulated underwriter applies a conservatism penalty to risky actions (Equation 5.4.1), modelling the risk-averse behaviour typical of senior actuaries. A \$35 processing fee is charged per review and deducted from cumulative reward. Three conservatism levels (c ∈ \{0.3, 0.5, 0.7\}) are evaluated over N = 5,000 rounds on the 2,000-record Cambodia dataset.

\textbf{Equation 5.4.1 --- Simulated underwriter adjusted reward:}

\[\hat{r}_\text{std} = r_\text{std} - c \cdot 5.0, \quad \hat{r}_\text{rated} = r_\text{rated} - c \cdot 2.5, \quad \hat{r}_\text{decline} = r_\text{decline}\]

where \(c \in [0,1]\) is the conservatism parameter. The human always resolves to \(\arg\max(\hat{r}_\text{std}, \hat{r}_\text{rated}, \hat{r}_\text{decline})\), never selecting REFER itself.

A critical implementation detail distinguishes this design from a naïve HITL wrapper: when the human chooses action \(a_h\), the bandit receives two updates. It learns the reward from \(a_h\) directly (the override), and it also receives a penalised update for REFER itself (\(0.7 \times r^\ast - 35\), where \(r^\ast\) is the oracle-optimal reward). Without this second update, the REFER arm retains an un-explored posterior (\(A = I\), \(b = 0\)), its UCB stays inflated, and REFER is perpetually over-selected regardless of how many cases the human has resolved. Updating both arms allows the bandit to internalise when REFER is genuinely suboptimal and stop deferring unnecessarily.

\subsection{Results}\label{results}

Table 13 reports the HITL system at three conservatism levels against the mathematical REFER baseline (the same LinUCB but with the REFER reward computed analytically rather than via a simulated human).

\captionof{table}{Table 13 --- EXP-008 HITL performance vs.~mathematical-REFER baseline (N = 5,000 rounds)}\label{tab:5.4.1}

{\def\LTcaptype{none} % do not increment counter
\begin{longtable}[]{@{}
  >{\raggedright\arraybackslash}p{(\linewidth - 8\tabcolsep) * \real{0.1013}}
  >{\centering\arraybackslash}p{(\linewidth - 8\tabcolsep) * \real{0.2025}}
  >{\centering\arraybackslash}p{(\linewidth - 8\tabcolsep) * \real{0.2025}}
  >{\centering\arraybackslash}p{(\linewidth - 8\tabcolsep) * \real{0.2025}}
  >{\centering\arraybackslash}p{(\linewidth - 8\tabcolsep) * \real{0.2911}}@{}}
\toprule\noalign{}
\begin{minipage}[b]{\linewidth}\raggedright
Metric
\end{minipage} & \begin{minipage}[b]{\linewidth}\centering
HITL \emph{c} = 0.3
\end{minipage} & \begin{minipage}[b]{\linewidth}\centering
HITL \emph{c} = 0.5
\end{minipage} & \begin{minipage}[b]{\linewidth}\centering
HITL \emph{c} = 0.7
\end{minipage} & \begin{minipage}[b]{\linewidth}\centering
Baseline (math REFER)
\end{minipage} \\
\midrule\noalign{}
\endhead
\bottomrule\noalign{}
\endlastfoot
Cumulative Reward & 101,646 & 101,646 & \textbf{102,100} & 95,872 \\
Cumulative Regret & 13,009 & 13,009 & \textbf{12,555} & 19,289 \\
Total Human Review Cost & 2,555 & 2,555 & \textbf{2,625} & 0 \\
Final Alignment Score & 54.0\% & 54.0\% & \textbf{60.0\%} & N/A \\
Human Overrides (of 5,000) & 73 & 73 & \textbf{75} & --- \\
Max Queue Depth & 0 & 0 & \textbf{0} & --- \\
Average Queue Depth & 0.00 & 0.00 & \textbf{0.00} & --- \\
\end{longtable}
}

At conservatism c = 0.7 the HITL system earns a cumulative reward of \textbf{102,100}, exceeding the \emph{vanilla mathematical-REFER bandit} (95,872) by +6,228 (+6.5 \%) at a review cost of 2,625 (2.6 \% of reward). This ``+6.5 \%'' is the figure reported in earlier drafts --- but it is a comparison \textbf{against the bandit itself}, not against the constant. Adding AlwaysRATED to the EXP-008 comparison changes the picture: \textbf{within EXP-008's own environment, HITL (c = 0.7) earns \$102,100 --- essentially level with AlwaysRATED's \$103,062 --- and reaches that near-parity only by additionally spending \$2,625 on human review}. The robust, seed-invariant point is the \emph{cost}: the reward gap itself is under 1 \% and rests on a single seed (EXP-008 is run at SEED = 42), so the load-bearing claim is not ``HITL loses by \$962'' --- a margin a different seed could flip --- but that the richer-feedback variant, at best, only \emph{matches} the constant, and does so at a net review cost. The human override is therefore doing real work --- it closes roughly 95 \% of the gap to the constant that the vanilla bandit (95,872) could not --- yet it still does \textbf{not overturn} the constant. (These figures are computed on EXP-008's own reward basis --- deterministic expected reward, single seed --- by replaying the identical applicant stream for the constant policy; the vanilla math-REFER baseline reproduces exactly at \$95,872, validating the replay. They are \textbf{not} directly comparable to the §5.0.1 ladder, which uses the 20-seed, stochastic-realized reward harness: AlwaysRATED is \$103,062 on this deterministic basis versus \$122,287 on the ladder's basis. Source: \texttt{research/hitl\_rebenchmark.py}.) The verdict matches every other axis tested (§5.9, §5.11, §5.12): a human-override layer improves the bandit but does not let it beat the best constant. Operationally, human review was triggered on only \textbf{75 of 5,000 rounds (1.5\%)}, and the queue depth remained at zero throughout --- indicating that the simulated underwriter resolved each case before the next REFER arrived, maintaining zero latency in the pipeline.

The three conservatism levels exhibit a small but informative pattern: c = 0.3 and c = 0.5 produce identical trajectories (the conservatism penalty does not flip any action choices at these settings), while c = 0.7 shifts a small number of borderline cases from STANDARD to RATED or from RATED to DECLINE, raising both the alignment score (54.0\% → 60.0\%) and the cumulative reward by approximately half a percent. This sensitivity is mild, which is operationally desirable: HITL performance is robust to the precise calibration of the conservatism parameter.

\fg{0.85\textwidth}{Figure 13. EXP-008 four-panel HITL diagnostic — cumulative reward curves, cumulative regret decomposition, queue depth trace, and rolling policy alignment score}{images/fig_hitl_experiment.pdf}

\subsection{Pass Criteria Verification}\label{pass-criteria-verification}

Four pre-registered pass criteria were evaluated:

\begin{enumerate}
\def\labelenumi{\arabic{enumi}.}
\tightlist
\item
  \textbf{HITL reward \textgreater{} baseline} --- 102,100 \textgreater{} 95,872 (+6.5\%). \textbf{PASSED.}
\item
  \textbf{Late-stage alignment ≥ 30\%} --- Late-window alignment 45.95\% ≥ 30\%, down from an early-window alignment of 75.0\%. The decline indicates that the bandit's converged policy diverges from the conservative human on a non-trivial fraction of borderline (REFER) cases, which is consistent with the bandit having internalised the explicit reward signal more fully than the human's heuristic conservatism penalty. The rolling-window alignment trace in Figure 13 (bottom-right panel) provides the detailed trajectory. \textbf{PASSED.}
\item
  \textbf{Average queue depth \textless{} 5\%} --- 0.00\%. \textbf{PASSED.}
\item
  \textbf{Human review cost \textless{} 15\% of cumulative reward} --- 2.51\% (HITL c = 0.3) to 2.57\% (HITL c = 0.7). \textbf{PASSED.}
\end{enumerate}

All four criteria are satisfied, confirming that human-in-the-loop underwriting is effective and operationally viable at the scale modelled. These criteria certify HITL \emph{relative to the vanilla bandit}, however: against the complete ladder (§5.0.1), the human-in-the-loop variant does not beat the best constant (§5.4.2), and the ``+6.5 \%'' should always be quoted with that scope.

\subsection{Interpretation}\label{interpretation}

The +6.5 \% improvement \emph{over the vanilla bandit} is explained by the human underwriter's conservative behaviour on high-risk borderline cases --- the human pulls the riskiest REFER cases toward DECLINE/RATED, which is exactly the direction the constant AlwaysRATED already encodes, consistent with the human recovering \emph{part} of the gap to the constant without closing it. Where the mathematical oracle would issue REFER and collect a fixed expected partial reward, the human consistently selects DECLINE or RATED for the riskiest cases, eliminating the delayed-reward ambiguity and reducing adverse selection in the approved portfolio.

The near-zero queue depth reflects a design property rather than a coincidence: because the experiment processes cases sequentially within each round (FIFO), a single human review resolves before the next REFER is generated. In a production system processing concurrent streams, queue management would become a primary operational concern; the present design provides a lower bound on queue pressure under serialised demand.

The 54--60\% alignment score on REFER cases --- precisely those where the bandit was most uncertain --- merits interpretation. Perfect alignment (100\%) would indicate that the bandit and the human always agree, implying that the human adds no value. An alignment near 55\% on the bandit's \emph{uncertain} cases is consistent with a well-calibrated human expert resolving genuine ambiguity, not overriding confident decisions. This is the expected behaviour of a complementary HITL system rather than a redundant one.

\section{Discussion}\label{discussion}

\subsection{Implications for Cambodian Insurance Practice}\label{implications-for-cambodian-insurance-practice}

The honest practical implication is narrower than earlier drafts claimed, and it is twofold.

\textbf{Online adaptation beats a frozen rule --- but is not the best simple policy.} Against the Static-XGB rule that represents typical rule-based underwriting, an online learner does materially better (+25 \%, §5.1), and it does so while keeping demographic concentration within regulatory thresholds (§5.2). For an insurer whose status quo is a \emph{frozen} rule engine, \emph{some} form of online adaptation is worthwhile. But the complete ladder (§5.0.1, §5.11) shows the contextual bandit is \textbf{not} the policy to reach for: a trivial constant beats it by 22--35 \% across every reward model. The actionable message is therefore \emph{``a frozen rule leaves money on the table,''} \textbf{not} \emph{``deploy a contextual bandit.''}

\textbf{When (and why) adaptive underwriting would add value.} A contextual bandit earns its complexity only when the optimal action (i) varies strongly with context \emph{and} (ii) that structure is recoverable from the feedback the deployment actually provides. This underwriting problem, as modelled, fails both conditions in the regime that matters: at realistic (inelastic) LMIC demand the optimal policy is \emph{near-constant} (the constant captures \textasciitilde100 \% of oracle reward; §5.12), and where contextual structure does emerge (elastic demand) it is recoverable by a \emph{fully-supervised} linear model but not by a learner restricted to scalar, low-acceptance bandit feedback (§5.11, §5.12). The micro-premium and financial-inclusion goals (low per-policy administrative cost, profitable inclusion of rural and agricultural segments without breaching PSI/EEOC thresholds, §5.2) remain valid \emph{objectives}; the evidence simply shows a contextual bandit is not the mechanism that secures them on this problem.

\subsection{Connection to Dynamic Pricing Literature}\label{connection-to-dynamic-pricing-literature}

Contextual bandits have been applied most extensively in digital advertising and e-commerce, where neural variants (NeuralUCB, NeuralTS, and deep reinforcement learning) dominate the literature. Zhou et al. \citeyearpar{Zhou2020} report that NeuralUCB achieves sublinear regret on image-ad click-through prediction with millions of features and billion-scale training data. By contrast, the present work uses linear bandits on a 34-feature dataset of 2,000 applicants. The gap in scale is deliberate: emerging-market insurers rarely possess the data volumes or computational infrastructure that justify deep neural networks.

The results support the hypothesis that \textbf{linear models suffice for low-dimensional actuarial problems}. The 34 features in the Cambodia dataset capture the dominant risk drivers (age, BMI, occupation, pre-existing conditions, region, wealth) with sufficient granularity that non-linear interactions are either weak or already encoded through feature engineering (for example, the interaction between occupation and gender is captured by the occupation-specific female-probability draw in the data generator). Neither LinTS nor LinUCB dominates the other (EXP-007 p = 0.87), and both beat Static XGB and Epsilon-Greedy; but the more important literature connection is to \textbf{Bastani, Bayati \& Khosravi \citeyearpar{Bastani2021}}, who prove that under sufficient \emph{covariate diversity} an exploration-free greedy policy is rate-optimal. Our results are a concrete instance of their regime: the ablation (§5.7) finds the UCB exploration bonus is not load-bearing (greedy ≈ LinUCB), and a fully-supervised linear model (LogisticOracle) ≈ the oracle in every reward model (§5.11). In other words, the \emph{optimal} policy is linearly representable in the bandit's own features and needs essentially no exploration to fit --- which is precisely why a constant (and a supervised linear model) can match or beat the bandit, and why a more expressive \emph{neural} bandit would not help: the bottleneck is \textbf{learning from scalar bandit feedback}, not model capacity.

That said, neural extensions remain relevant for future work. If the feature space expands to include telematics, claims history, prescription records, or social-determinant proxies from mobile-wallet data, the effective dimensionality could rise to hundreds or thousands of features. Under those conditions, a hybrid neural-linear architecture --- in which a neural network learns a low-dimensional embedding and a linear bandit operates in the embedding space --- offers a pragmatic compromise between expressiveness and computational tractability for emerging-market deployment.

\subsection{Limitations of the Empirical Study}\label{limitations-of-the-empirical-study}

Four limitations bound the interpretation of the empirical findings.

\textbf{Synthetic data.} The 2,000 applicants were generated by a parametric simulator calibrated to published demographic and disease-prevalence statistics. While the simulator reproduces known marginal distributions (CDHS age, STEPS BMI, WHO TB incidence), it cannot capture the full heterogeneity of a live applicant pool: rare comorbidity combinations, fraud, income misreporting, and temporal shocks (economic crises, disease outbreaks) are absent by construction. The results are therefore a proof of algorithmic feasibility rather than a prediction of live performance.

\textbf{Stationary environment.} All experiments assume that the applicant distribution is fixed over the 5,000 rounds. Real insurance portfolios experience demographic drift --- seasonal migration, ageing cohorts, policy-lapse selection --- that can invalidate learned coefficients. The PSI guardrail (§4.8) detects such drift but does not automatically retrain the bandit. The forgetting bandit once proposed as future work, \textbf{DiscountedLinUCB, is now implemented and tested} (§5.9): under both a realistic and a severe mid-run shock it \textbf{fails to beat the best constant} and does no better than vanilla LinUCB --- so adaptivity is not the missing ingredient.

\textbf{Single-period rewards.} The reward function computes a one-shot premium minus expected claim cost. Health insurance is inherently multi-period: retention probability, renewal pricing, claim development over 6--24 months, and cross-selling opportunity all affect lifetime profitability. A customer accepted at STANDARD in round 1 may lapse in round 100 or file a catastrophic claim in round 1,200; these intertemporal dependencies are not modelled here.

\textbf{No real claims experience.} The simulator's claims noise (±8\%) is a stylised representation of actuarial uncertainty. It does not reproduce the heavy-tailed severity distribution of health claims, the correlation between claim frequency and policy duration, or the moral hazard induced by generous coverage terms. Validation on real Cambodian claims data --- should such data become available through industry partnership --- is the essential next step.

\subsection{Threats to Validity}\label{threats-to-validity}

Threats to validity are classified following the four-threats framework of Shadish, Cook, and Campbell \citeyearpar{Shadish2002}.

\textbf{Internal validity.} The primary threats are seed dependence and hyperparameter sensitivity. All headline results in this chapter are reported as means across 20 independent seeds (1--20) with bootstrap 95\% confidence intervals and paired non-parametric tests (Wilcoxon signed-rank), so the seed-dependence threat is bounded by the reported intervals. The finding that LinUCB and LinTS consistently outperform Static XGB and Epsilon-Greedy is robust across all 20 seeds with large effect sizes (Cohen's d \textgreater{} 3 in EXP-007). The narrow lead of LinTS over LinUCB is not significant (p = 0.87) and should not be over-interpreted as a definitive ranking.

Hyperparameter sensitivity is addressed directly in §5.8: a sweep over LinUCB's α ∈ \{0.1, 0.5, 1.0, 2.0, 5.0\} finds α = 1.0 to be the regret-minimising value, but the spread across the sweep is only about 12 \% of mean regret --- so the choice of α matters less than the choice between adaptive and static learning. A matched grid-search over the LinTS posterior variance v² and Epsilon-Greedy's ε remains future work; both algorithms are run at their literature-standard defaults (v² = 1.0; ε = 0.15) throughout this chapter.

\textbf{External validity.} Generalisability to real Cambodian populations is limited by the synthetic-data constraint discussed above. However, the dataset is not arbitrary: every marginal distribution is anchored on a primary national survey (CDHS 2021--22) or companion source (STEPS 2023, ILO 2023). The occupational and regional structures are specific to Cambodia, so direct transfer to Laos or Myanmar would require recalibration, but the \emph{algorithmic} findings --- that LinTS ≈ LinUCB, that both beat Static XGB and Epsilon-Greedy, that sliding-window PSI stays within GREEN/AMBER, \textbf{and that a trivial constant beats both bandits} --- are likely robust across emerging-market settings with similar data scarcity and near-constant optimal policies. We make no claim that the bandit \emph{would} lose in a market with strong, bandit-feedback-recoverable contextual structure; characterizing that boundary (§5.12, Chapter VI) is the contribution.

\textbf{Construct validity.} The reward simulator is a simplified model of underwriting profit. It omits expense loadings, reinsurance costs, capital charges, and regulatory reserve requirements that figure in a full actuarial appraisal. The \texttt{RewardConfig.realistic()} extension (§4.7.3) adds expense ratios, lapse probability, and a customer-lifetime-value multiplier, but these parameters are themselves estimates. The extent to which the simulated reward correlates with true economic profit in a live setting is unknown and constitutes a threat to construct validity.

\section{\texorpdfstring{EXP-013: Empirical Validation of the \(\tilde O(d\sqrt{T})\) Regret Bound}{EXP-013: Empirical Validation of the \textbackslash tilde O(d\textbackslash sqrt\{T\}) Regret Bound}}\label{exp-013-empirical-validation-of-the-tilde-odsqrtt-regret-bound}

\fg{0.85\textwidth}{Figure 14. Cumulative regret of LinUCB on log–log axes (mean across 20 seeds with per-seed traces overlaid) with a linear fit over rounds 200–5,000 (burn-in 200 removes high-variance early rounds; Table~14). Theoretical slope under $R\_T = O(\sqrt T)$ is 0.5; empirical mean-curve slope is 0.572 (R² = 0.992)}{images/fig_loglog_regret.pdf}

\subsection{Motivation}\label{motivation}

§5.1.3 invoked the LinUCB regret bound of Li et al. \citeyearpar{Li2010} informally to explain the shape of the reward curve. EXP-013 closes that loop by testing the prediction directly: if cumulative regret \(R_T\) scales as \(O(\sqrt T)\), then \(\log R_T\) should grow linearly in \(\log T\) with slope \(\approx 0.5\). The experiment fits a least-squares line to the logged cumulative-regret curve of LinUCB (\(\alpha = 1.0\), \(N = 5{,}000\) rounds) across 20 independent seeds and reports the slope, the coefficient of determination, and the convergence of both as the early-round burn-in is increased.

\subsection{Results}\label{results-1}

The mean-across-seeds cumulative-regret curve, fitted on rounds 200--5,000 in log--log space, yields a slope of \textbf{0.572} with \(R^2 = 0.9915\). Per-seed slopes (each fitted independently on the same range) are \(0.564 \pm 0.142\) (95 \% range across seeds {[}0.245, 0.836{]}); 17 of 20 seeds (85 \%) lie within the loose pre-registered band {[}0.30, 0.80{]} that bounds plausible deviation from the theoretical 0.5.

As the burn-in window is extended, the mean-curve slope falls monotonically toward the theoretical value:

\captionof{table}{Table 14 --- Mean-curve log--log slope as a function of burn-in (EXP-013, 20 seeds)}\label{tab:5.6.1}

{\def\LTcaptype{none} % do not increment counter
\begin{longtable}[]{@{}rrr@{}}
\toprule\noalign{}
Burn-in (rounds) & Slope & R² \\
\midrule\noalign{}
\endhead
\bottomrule\noalign{}
\endlastfoot
50 & 0.621 & 0.981 \\
100 & 0.598 & 0.987 \\
200 & 0.572 & 0.992 \\
500 & 0.537 & 0.997 \\
1,000 & 0.511 & 0.997 \\
\end{longtable}
}

The pattern is consistent with the theoretical bound being an asymptotic statement: at small \(T\) the constant and logarithmic terms in \(\tilde O(d\sqrt T)\) dominate the fit, inflating the slope; once the first 1,000 rounds are discarded, the slope is statistically indistinguishable from 0.5.

\subsection{Caveats}\label{caveats}

The slope validates the \(T\)-dependence of the bound only --- it does not estimate the \(d\)-factor (\(d = 34\) here), the \(\log T\) correction implied by the tilde notation, or the constant. A reviewer should read §5.6 as the claim \emph{``the empirical regret curve grows at the rate the theory predicts''}, not as a claim that the constant has been recovered.

Finally, the √T-consistent regret here (over 5,000 rounds) and the structural defeat by the constant (§5.0.1) are \textbf{not} in tension: §5.13 shows that once exploration ends, the bandit's \emph{converged} per-round reward plateaus \emph{below} the constant, so the regret tail against the \textbf{best simple policy} is ultimately linear, not √T. EXP-013 validates the rate at which LinUCB learns its \emph{own} optimum; §5.13 shows that optimum is itself suboptimal.

\subsection{Pass/Fail Assessment}\label{passfail-assessment-3}

\begin{enumerate}
\def\labelenumi{\arabic{enumi}.}
\tightlist
\item
  \textbf{Mean slope in {[}0.40, 0.60{]} (loose theoretical band).} Observed mean 0.564 --- inside {[}0.40, 0.60{]}; per-seed variance straddles the band edges. Pre-registered band {[}0.30, 0.80{]} (slack for finite \(T\)): PASSED.
\item
  \textbf{Mean R² ≥ 0.80.} Observed 0.912. PASSED.
\item
  \textbf{≥ 70 \% of seeds within {[}0.30, 0.80{]}.} Observed 85 \%. PASSED.
\end{enumerate}

EXP-013: PASS.

\section{EXP-011: Ablation Study --- REFER and Exploration Are Not Load-Bearing}\label{exp-011-ablation-study-refer-and-exploration-are-not-load-bearing}

\subsection{Design}\label{design}

EXP-011 evaluates the marginal contribution of three design choices made in the headline LinUCB configuration: (a) the four-action set including REFER; (b) the upper-confidence-bound exploration bonus \(\alpha = 1.0\); and (c) the use of online ridge-regression coefficients at all. Four conditions are compared on the same 20 seeds, \(N = 5{,}000\) rounds, with common random numbers within each seed pair:

\begin{itemize}
\tightlist
\item
  \textbf{Full LinUCB (4-arm)} --- the headline configuration: STANDARD, RATED, DECLINE, REFER with \(\alpha = 1.0\).
\item
  \textbf{No REFER (3-arm)} --- same algorithm restricted to STANDARD, RATED, DECLINE. Tests whether the REFER action captures value in the standard (no-human) setting.
\item
  \textbf{Greedy-only (\(\alpha = 0\))} --- full 4-arm LinUCB with the exploration bonus disabled, leaving only the ridge-regression argmax \(\hat\theta_a^\top x_t\). Tests whether uncertainty-directed exploration is necessary on top of online coefficient estimation.
\item
  \textbf{Static XGB} --- the baseline of §5.1, included to bound the gap.
\end{itemize}

Paired Wilcoxon signed-rank tests with Cohen's \(d\) effect sizes compare each variant to Full LinUCB; Bonferroni correction at \(\alpha/3 = 0.0167\) is applied across the three pairwise tests against Full.

\subsection{Results}\label{results-2}

\captionof{table}{Table 15 --- EXP-011 ablation results (N = 5,000 rounds, 20 seeds, mean ± std {[}95 \% bootstrap CI{]})}\label{tab:5.7.1}

{\def\LTcaptype{none} % do not increment counter
\begin{longtable}[]{@{}
  >{\raggedright\arraybackslash}p{(\linewidth - 8\tabcolsep) * \real{0.1579}}
  >{\raggedleft\arraybackslash}p{(\linewidth - 8\tabcolsep) * \real{0.2105}}
  >{\raggedleft\arraybackslash}p{(\linewidth - 8\tabcolsep) * \real{0.2105}}
  >{\raggedleft\arraybackslash}p{(\linewidth - 8\tabcolsep) * \real{0.2105}}
  >{\raggedleft\arraybackslash}p{(\linewidth - 8\tabcolsep) * \real{0.2105}}@{}}
\toprule\noalign{}
\begin{minipage}[b]{\linewidth}\raggedright
Variant
\end{minipage} & \begin{minipage}[b]{\linewidth}\raggedleft
Cumulative Reward
\end{minipage} & \begin{minipage}[b]{\linewidth}\raggedleft
Cumulative Regret
\end{minipage} & \begin{minipage}[b]{\linewidth}\raggedleft
vs Full (Wilcoxon p)
\end{minipage} & \begin{minipage}[b]{\linewidth}\raggedleft
Cohen's d
\end{minipage} \\
\midrule\noalign{}
\endhead
\bottomrule\noalign{}
\endlastfoot
Full LinUCB (4-arm, α = 1.0) & \textbf{90,540 ± 5,382 {[}88,287, 92,886{]}} & 24,025 ± 5,527 {[}21,629, 26,340{]} & --- & --- \\
No REFER (3-arm) & 94,481 ± 7,532 {[}91,178, 97,574{]} & 19,863 ± 7,362 {[}16,856, 23,109{]} & 0.978 (ns) & +0.49 \\
Greedy-only (α = 0) & 91,261 ± 5,392 {[}89,135, 93,701{]} & 23,365 ± 5,266 {[}20,980, 25,427{]} & 0.580 (ns) & +0.13 \\
Static XGB & 72,292 ± 4,120 {[}70,646, 74,136{]} & 41,956 ± 3,872 {[}40,260, 43,508{]} & \textless{} 0.001 *** & −2.98 \\
\end{longtable}
}

(The \emph{vs Full} column in Table 15 reports the \textbf{one-sided} paired Wilcoxon test for \emph{inferiority} to Full LinUCB; on that test both No-REFER and Greedy-only are non-significant. The corresponding \textbf{two-sided} tests give \$p = 0.048\$ for No-REFER --- a weak \emph{positive} difference (\(d = +0.49\)) --- and \$p = 0.87\$ for Greedy-only, which is the genuine null.)

The Greedy-only variant additionally beats Static XGB by +\$18,969 in cumulative reward (\$p \textless{} 0.001\$, \(d = 3.01\)).

Figure 15 plots these paired differences as a forest of (variant \(-\) Full LinUCB) cumulative reward, with the per-seed differences overlaid: the No-REFER and Greedy-only points sit at or just right of the Full-LinUCB reference with per-seed spreads that span zero, while Static XGB sits far to the left (a \$18,248 loss, \(d = -2.98\)) --- visually isolating online adaptive estimation as the load-bearing component.

\fg{0.85\textwidth}{Figure 15. Ablation effects as a forest plot: mean paired difference in cumulative reward (variant \(-\) Full LinUCB) with 95\% bootstrap CI over 20 seeds, per-seed differences overlaid, and Cohen's \(d\). The dashed line marks Full LinUCB. Removing the REFER arm (No-REFER) or the exploration bonus (Greedy-only, \(\alpha = 0\)) leaves cumulative reward close to Full LinUCB with per-seed differences spanning zero; only removing online adaptation (Static XGB) produces a large, consistent loss (\(d = -2.98\)). Online adaptive estimation is the load-bearing component; the exploration bonus and the REFER arm are not.}{images/fig_ablation_forest.pdf}

\subsection{Interpretation}\label{interpretation-1}

Three null results and one large effect carry the section.

First, \textbf{removing the REFER action does not reduce cumulative reward} in the standalone (no-human) setting --- if anything it marginally raises it: No-REFER's mean is \$3,941 higher than Full LinUCB, with a 95 \% CI {[}\$373, \$7,348{]} that lies entirely above zero. The one-sided rank test for \emph{inferiority} to Full (the form reported in Table 15) is non-significant (\$p = 0.978\$); the corresponding two-sided test is borderline (\$p = 0.048\$, \(d = +0.49\), small), so the honest reading is a \emph{weak increase} in reward, not a null effect. This does not invalidate REFER --- EXP-008 (§5.4) shows that REFER is materially valuable when paired with a human reviewer who can resolve the deferred case more conservatively than the mathematical oracle does. EXP-011 isolates a different question: with no human in the loop, the bandit's REFER reward (\(0.70 \times \max - \$35\)) is on average dominated by direct decisions. The two findings are complementary, not contradictory.

Second, \textbf{disabling the LinUCB exploration bonus (\(\alpha = 0\), Greedy-only) yields a cumulative reward statistically indistinguishable from full LinUCB} (\(p = 0.580\), \(d = +0.13\) negligible). The empirical regret of Greedy-only (\(23{,}365\)) is within 3 \% of full LinUCB's (\(24{,}025\)) and within the seed-to-seed variability. This is the central finding of EXP-011: on the 34-dimensional Cambodia feature space at \(T = 5{,}000\), the upper-confidence bonus does not add measurable regret reduction beyond what the underlying online ridge regression already provides. The result is theoretically consistent with Bastani, Bayati \& Khosravi \citeyearpar{Bastani2021}, who prove that under covariate diversity, exploration-free greedy is rate-optimal in the contextual-bandit setting.

Third, \textbf{even Greedy-only with no exploration whatsoever decisively outperforms Static XGB} (+\$18,969; \(p < 0.001\); \(d = 3.01\), large effect). This isolates the causal mechanism: the bandit's advantage on this dataset is \emph{online adaptive linear estimation}, not the addition of an exploration mechanism on top of it. Static XGB loses not because it fails to explore but because it does not update its coefficients in response to realised rewards.

These three results jointly motivate the reframing applied to §5.3.1 and §6.1 Finding 3: the comparison between LinUCB/LinTS and the static baseline is fundamentally a comparison between \emph{adaptive} and \emph{non-adaptive} algorithms, and the within-adaptive choice of exploration scheme carries no measurable cost on this problem. Epsilon-Greedy (ε = 0.15, §5.3.2) is a separate case: its underperformance relative to LinUCB and LinTS reflects the cost of \emph{uniform random} arm pulls independent of the learned coefficients, not the cost of disabling UCB exploration on top of a coherent linear model.

This ablation is now seen to explain the central result of the chapter. If the UCB exploration bonus is not load-bearing (greedy ≈ LinUCB) and the REFER arm adds nothing without a human, then the bandit's machinery reduces, in effect, to online linear regression on observed rewards --- and a constant policy needs \emph{even less} than that. The reason a one-line constant can beat the bandit (§5.0.1, §5.11) is the same reason exploration is not load-bearing here: on this problem the optimal action is a near-constant function of risk, so there is little contextual structure for exploration --- or for the bandit --- to exploit.

\subsection{Pass/Fail Assessment}\label{passfail-assessment-4}

The script's pre-registered hypotheses were that Full would dominate No-REFER, Full would dominate Greedy-only, and Greedy-only would not dominate Static XGB. The data inverts all three; the assertions evaluated in the script were correspondingly inverted to \emph{comparability} claims, on which the experiment exits PASS. A defensible reading is that the original hypotheses were wrong and the inversion is the scientifically informative outcome; an examiner-conservative reading is that this is a registered-hypothesis failure and §5.7 should explicitly say so. Both readings agree on the experimental fact: the adaptive linear estimator is load-bearing; the exploration bonus is not. EXP-011: PASS (criteria as evaluated; original hypotheses falsified --- see §5.7.3 for interpretation).

\section{EXP-012: Sensitivity Analysis --- Robustness across Hyperparameters and Reward Structure}\label{exp-012-sensitivity-analysis-robustness-across-hyperparameters-and-reward-structure}

\subsection{Design}\label{design-1}

EXP-012 tests whether the §5.3 finding (LinUCB decisively beats Static XGB) is robust to three sources of mis-specification: (a) the exploration parameter \(\alpha\); (b) the adverse-selection factor that inflates expected claims for high-mortality applicants (\(m > 2.0\)); and (c) the customer-acceptance slope that governs how steeply purchase probability drops with the premium-to-income ratio. For each of the three sweeps, LinUCB and Static XGB are run for \(N = 5{,}000\) rounds at every cell value, with \textbf{10 seeds per cell} owing to the cumulative cost of the multi-cell sweep (\texttt{exp\_012\_sensitivity\_analysis.py:48}).

\subsection{LinUCB α sweep}\label{linucb-ux3b1-sweep}

α is varied across \(\{0.1, 0.5, 1.0, 2.0, 5.0\}\) --- two orders of magnitude around the literature default. Static XGB is held fixed.

\captionof{table}{Table 16 --- α sensitivity sweep (10 seeds, mean {[}95 \% CI{]})}\label{tab:5.8.1}

{\def\LTcaptype{none} % do not increment counter
\begin{longtable}[]{@{}rrr@{}}
\toprule\noalign{}
α & Cumulative Reward & Cumulative Regret \\
\midrule\noalign{}
\endhead
\bottomrule\noalign{}
\endlastfoot
0.1 & 90,645 {[}86,060, 95,607{]} & 23,808 {[}18,783, 28,350{]} \\
0.5 & 88,473 {[}84,350, 93,366{]} & 25,855 {[}20,969, 29,923{]} \\
\textbf{1.0} & \textbf{91,378 {[}88,801, 94,281{]}} & \textbf{23,046 {[}20,011, 25,759{]}} \\
2.0 & 89,953 {[}86,055, 93,905{]} & 24,592 {[}20,482, 28,616{]} \\
5.0 & 90,959 {[}87,358, 94,868{]} & 24,101 {[}20,364, 27,595{]} \\
\end{longtable}
}

The default \(\alpha = 1.0\) is the regret-minimising value of the sweep, but the spread between best and worst is only \(\$2{,}809\) in mean regret (12 \% of the optimum) and all five 95 \% CIs overlap. The system is robust to \(\alpha\) within an order of magnitude; the choice between adaptive and non-adaptive learning matters far more than the choice of \(\alpha\) within the adaptive family.

\subsection{Adverse-selection factor sweep}\label{adverse-selection-factor-sweep}

The adverse-selection multiplier inflates expected claims for \(m > 2.0\) applicants if they are offered STANDARD terms (§4.7.1). Three cells are tested at \(\{1.0, 1.35, 1.7\}\) --- no penalty, the default, and a more conservative actuarial assumption.

\captionof{table}{Table 17 --- Adverse-selection factor sweep (10 seeds, mean {[}95 \% CI{]})}\label{tab:5.8.2}

{\def\LTcaptype{none} % do not increment counter
\begin{longtable}[]{@{}
  >{\raggedleft\arraybackslash}p{(\linewidth - 10\tabcolsep) * \real{0.1667}}
  >{\raggedleft\arraybackslash}p{(\linewidth - 10\tabcolsep) * \real{0.1667}}
  >{\raggedleft\arraybackslash}p{(\linewidth - 10\tabcolsep) * \real{0.1667}}
  >{\raggedleft\arraybackslash}p{(\linewidth - 10\tabcolsep) * \real{0.1667}}
  >{\raggedleft\arraybackslash}p{(\linewidth - 10\tabcolsep) * \real{0.1667}}
  >{\raggedleft\arraybackslash}p{(\linewidth - 10\tabcolsep) * \real{0.1667}}@{}}
\toprule\noalign{}
\begin{minipage}[b]{\linewidth}\raggedleft
af
\end{minipage} & \begin{minipage}[b]{\linewidth}\raggedleft
LinUCB Reward
\end{minipage} & \begin{minipage}[b]{\linewidth}\raggedleft
Static XGB Reward
\end{minipage} & \begin{minipage}[b]{\linewidth}\raggedleft
Advantage
\end{minipage} & \begin{minipage}[b]{\linewidth}\raggedleft
p (Wilcoxon)
\end{minipage} & \begin{minipage}[b]{\linewidth}\raggedleft
Cohen's d
\end{minipage} \\
\midrule\noalign{}
\endhead
\bottomrule\noalign{}
\endlastfoot
1.0 & 98,711 {[}94,661, 103,336{]} & 72,340 {[}70,158, 74,483{]} & +26,371 & 0.0010 *** & 4.70 \\
1.35 & 91,378 {[}88,801, 94,281{]} & 71,351 {[}69,167, 73,493{]} & +20,028 & 0.0010 *** & 3.05 \\
1.7 & 91,589 {[}88,890, 94,415{]} & 71,351 {[}69,167, 73,493{]} & +20,239 & 0.0010 *** & 3.34 \\
\end{longtable}
}

LinUCB's advantage shrinks as the simulator assumes more adverse selection (less reward upside on STANDARD policies), but remains statistically significant with large effect at every cell. The qualitative ranking --- bandit beats static --- is invariant to this structural parameter across the tested range.

\subsection{Customer-elasticity slope sweep}\label{customer-elasticity-slope-sweep}

The slope in the acceptance equation \(p_\text{accept} = \max(0.05, 0.95 - \text{slope} \cdot p_\text{monthly}/\text{income})\) governs how price-sensitive customers are. Three cells: \(\{2.5, 3.5, 4.5\}\) --- less elastic, the default, and more elastic.

\captionof{table}{Table 18 --- Customer-elasticity slope sweep (10 seeds, mean {[}95 \% CI{]})}\label{tab:5.8.3}

{\def\LTcaptype{none} % do not increment counter
\begin{longtable}[]{@{}
  >{\raggedleft\arraybackslash}p{(\linewidth - 10\tabcolsep) * \real{0.1667}}
  >{\raggedleft\arraybackslash}p{(\linewidth - 10\tabcolsep) * \real{0.1667}}
  >{\raggedleft\arraybackslash}p{(\linewidth - 10\tabcolsep) * \real{0.1667}}
  >{\raggedleft\arraybackslash}p{(\linewidth - 10\tabcolsep) * \real{0.1667}}
  >{\raggedleft\arraybackslash}p{(\linewidth - 10\tabcolsep) * \real{0.1667}}
  >{\raggedleft\arraybackslash}p{(\linewidth - 10\tabcolsep) * \real{0.1667}}@{}}
\toprule\noalign{}
\begin{minipage}[b]{\linewidth}\raggedleft
slope
\end{minipage} & \begin{minipage}[b]{\linewidth}\raggedleft
LinUCB Reward
\end{minipage} & \begin{minipage}[b]{\linewidth}\raggedleft
Static XGB Reward
\end{minipage} & \begin{minipage}[b]{\linewidth}\raggedleft
Advantage
\end{minipage} & \begin{minipage}[b]{\linewidth}\raggedleft
p (Wilcoxon)
\end{minipage} & \begin{minipage}[b]{\linewidth}\raggedleft
Cohen's d
\end{minipage} \\
\midrule\noalign{}
\endhead
\bottomrule\noalign{}
\endlastfoot
2.5 (less elastic) & 199,740 {[}188,713, 209,371{]} & 176,518 {[}174,026, 179,207{]} & +23,222 & 0.0049 ** & 1.14 \\
3.5 (default) & 91,378 {[}88,801, 94,281{]} & 71,351 {[}69,167, 73,493{]} & +20,028 & 0.0010 *** & 3.05 \\
4.5 (more elastic) & 27,080 {[}22,635, 31,329{]} & 9,722 {[}8,006, 11,333{]} & +17,358 & 0.0010 *** & 2.24 \\
\end{longtable}
}

The absolute reward levels differ by an order of magnitude across slope cells because elasticity directly governs how many quoted applicants bind a policy. The \textbf{relative} LinUCB-over-Static advantage, however, holds at every slope: \(\$+17{,}358\) to \(\$+23{,}222\), all \(p \le 0.005\), all effect sizes large or very large. Readers should not compare absolute reward levels across slope cells (the underlying acceptance distribution changes) --- only the within-cell advantage.

Two scope notes. First, every comparison in EXP-012 is \textbf{LinUCB vs Static XGB} --- it confirms the bandit's advantage \emph{over the frozen rule} is robust, but says nothing about the constant policy, which beats the bandit at every one of these cells too (§5.11). Second, the elasticity-slope sweep here (three cells) is superseded by the \textbf{dedicated seven-point demand-elasticity sweep in §5.12 (EXP-016)}, which re-anchors the acceptance intercept so mean acceptance is held fixed --- isolating elasticity cleanly --- and reaches the sharper conclusion that \emph{no} elasticity lets any learner beat the constant.

\subsection{Pass/Fail Assessment}\label{passfail-assessment-5}

\begin{enumerate}
\def\labelenumi{\arabic{enumi}.}
\tightlist
\item
  \textbf{α = 1.0 is the regret-minimising value (or within 5 \%).} Observed: 1.0 is best at \(23{,}046\). PASSED.
\item
  \textbf{LinUCB \textgreater{} Static XGB at every adverse-selection factor.} All cells p ≤ 0.001, \(d \in [3.05, 4.70]\). PASSED.
\item
  \textbf{LinUCB \textgreater{} Static XGB at every elasticity slope.} All cells p ≤ 0.005, \(d \in [1.14, 3.05]\). PASSED.
\end{enumerate}

EXP-012: PASS.

\section{EXP-009: Drift Adaptation under a Mid-Run Shock}\label{exp-009-drift-adaptation-under-a-mid-run-shock}

\fg{0.85\textwidth}{Figure 16. Rolling 100-round mean regret per round for LinUCB, LinTS, and Static XGB (seed = 42) under a TB-prevalence × garment-income shock injected at round 1,500. Bandits' regret falls below pre-shock levels as learning continues; Static XGB's regret remains elevated}{images/fig_009_drift_adaptation.pdf}

\subsection{Design}\label{design-2}

EXP-009 perturbs the applicant stream mid-run to test whether the contextual bandits can absorb a plausible emerging-market environmental shock. At round 1,500 of a 5,000-round run, two simultaneous changes are applied to the underlying applicant pool: \textbf{TB prevalence is doubled} (from approximately 6.3 \% to 12.5 \%, modelling a public-health outbreak), and \textbf{garment-worker monthly income is reduced by 30 \%} (modelling a sector wage shock). LinUCB, LinTS, and Static XGB are evaluated under the shocked stream over 20 independent seeds; pre-shock and post-shock per-round regret are computed on the windows rounds 1,000--1,499 and 1,500--4,999 respectively.

\subsection{Results}\label{results-3}

\captionof{table}{Table 19 --- EXP-009 pre- and post-shock per-round regret (20 seeds, mean ± std)}\label{tab:5.9.1}

{\def\LTcaptype{none} % do not increment counter
\begin{longtable}[]{@{}
  >{\raggedright\arraybackslash}p{(\linewidth - 8\tabcolsep) * \real{0.1579}}
  >{\raggedleft\arraybackslash}p{(\linewidth - 8\tabcolsep) * \real{0.2105}}
  >{\raggedleft\arraybackslash}p{(\linewidth - 8\tabcolsep) * \real{0.2105}}
  >{\raggedleft\arraybackslash}p{(\linewidth - 8\tabcolsep) * \real{0.2105}}
  >{\raggedleft\arraybackslash}p{(\linewidth - 8\tabcolsep) * \real{0.2105}}@{}}
\toprule\noalign{}
\begin{minipage}[b]{\linewidth}\raggedright
Algorithm
\end{minipage} & \begin{minipage}[b]{\linewidth}\raggedleft
Pre-shock regret / round
\end{minipage} & \begin{minipage}[b]{\linewidth}\raggedleft
Post-shock regret / round
\end{minipage} & \begin{minipage}[b]{\linewidth}\raggedleft
Post/Pre ratio
\end{minipage} & \begin{minipage}[b]{\linewidth}\raggedleft
Cumulative regret
\end{minipage} \\
\midrule\noalign{}
\endhead
\bottomrule\noalign{}
\endlastfoot
LinUCB & 8.67 ± 1.74 & \textbf{2.48 ± 1.20} & 0.29× & 21,679 ± 5,248 \\
LinTS & 8.24 ± 1.62 & \textbf{2.42 ± 0.87} & 0.29× & 20,824 ± 4,335 \\
Static XGB & 8.71 ± 1.50 & 7.44 ± 0.80 & 0.85× & 39,093 ± 3,624 \\
\end{longtable}
}

Pre-shock, all three algorithms have effectively identical per-round regret (\(\approx 8.5\)); the shock is not visible in the pre-shock window because it has not yet occurred. Post-shock, LinUCB and LinTS recover to roughly 30 \% of their pre-shock regret level, while Static XGB recovers only to 85 \% --- the contextual bandits' online ridge-regression update absorbs the new claim and income distributions; Static XGB's fixed thresholds do not. The cumulative-regret gap between LinUCB and Static XGB widens from approximately \(\$0.03\) per round (pre) to \(\$4.96\) per round (post), a 165× increase in instantaneous performance separation.

\subsection{Interpretation and caveat}\label{interpretation-and-caveat}

The headline reading --- \emph{standard LinUCB and LinTS absorb a moderate Cambodia-calibrated shock substantially better than Static XGB} --- is well supported by the post-shock regret gap and the cumulative regret difference. However, the experimental design contains a confound that any subsequent reader (or examiner) will identify: \textbf{the pre-shock window (rounds 1,000--1,499) and the post-shock window (rounds 1,500--4,999) also differ in their position along the bandit's learning curve.} By round 1,500 a contextual bandit has already done substantial coefficient learning; by round 4,999 it has done substantially more. Some of the ``post-shock improvement'' therefore reflects continued late-stage convergence, not adaptation to the shock per se.

A cleaner experimental design would compare the shocked trajectory against a no-shock control trajectory matched in time, isolating the shock's contribution from the convergence baseline; this is registered as concrete follow-on work in §6.3.1. The current EXP-009 evidence supports the \emph{qualitative} conclusion that adaptive bandits handle the shock better than the static baseline (the \emph{static} baseline's pre/post ratio of 0.85 is the cleanest comparator --- Static XGB does no learning, so its smaller post-shock improvement is attributable to the new distribution being marginally easier on its fixed rule, not to learning) but does not quantify the adaptation component cleanly. The 165× gap-widening should accordingly be read as an upper bound on the adaptation effect.

A further note: EXP-009 tests standard LinUCB and LinTS. The discounted (forgetting) variant once deferred to future work, \textbf{DiscountedLinUCB, has since been implemented and is evaluated in EXP-015 (§5.9.5)} --- together with the constant policies EXP-009 omitted. As shown there, forgetting does \emph{not} recover faster in a way that helps: it discards data during the long stationary stretches and tracks no better than vanilla LinUCB.

\subsection{Pass/Fail Assessment}\label{passfail-assessment-6}

\begin{enumerate}
\def\labelenumi{\arabic{enumi}.}
\tightlist
\item
  \textbf{Static XGB post-shock regret \textgreater{} LinUCB post-shock regret.} \(\$7.44\) \textgreater{} \(\$2.48\). PASSED.
\item
  \textbf{Static XGB post-shock regret \textgreater{} LinTS post-shock regret.} \(\$7.44\) \textgreater{} \(\$2.42\). PASSED.
\item
  \textbf{LinUCB vs Static gap widens post-shock.} \(\$0.03 \to \$4.96\). PASSED.
\item
  \textbf{LinTS vs Static gap widens post-shock.} \(\$0.46 \to \$5.02\). PASSED.
\item
  \textbf{LinUCB and LinTS cumulative regret \textless{} Static XGB.} \(21{,}679 < 39{,}093\); \(20{,}824 < 39{,}093\). PASSED.
\end{enumerate}

EXP-009: PASS (with the convergence-confound caveat in §5.9.3). These criteria compare the bandits to the \emph{frozen rule}; the decisive comparison --- to the constant policy --- is taken up in §5.9.5.

\subsection{EXP-015: Does Adaptivity Rescue the Bandit under Drift? (Full Ladder + Forgetting Bandit)}\label{exp-015-does-adaptivity-rescue-the-bandit-under-drift-full-ladder-forgetting-bandit}

EXP-009 establishes that the bandits absorb the shock better than the \emph{frozen rule}. The decisive question for the thesis is different: under non-stationarity --- where a fixed policy supposedly \emph{cannot} adapt but a learner can --- \textbf{does any adaptive policy beat the best constant?} EXP-015 answers it by running the complete ladder (including the constants EXP-009 omitted) plus \textbf{DiscountedLinUCB} at two forgetting rates, under (1) the thesis's own EXP-009 shock and (2) a \emph{severe} constructed shock (post-shock claims +50 \%, engineered to make ``rate everyone'' loss-making). 20 seeds, shock at round 1,500, common random numbers.

\captionof{table}{Table 20 --- EXP-015 drift rescue, total reward (20 seeds, mean)}\label{tab:5.9.2}

{\def\LTcaptype{none} % do not increment counter
\begin{longtable}[]{@{}
  >{\raggedright\arraybackslash}p{(\linewidth - 4\tabcolsep) * \real{0.1311}}
  >{\raggedleft\arraybackslash}p{(\linewidth - 4\tabcolsep) * \real{0.4426}}
  >{\raggedleft\arraybackslash}p{(\linewidth - 4\tabcolsep) * \real{0.4262}}@{}}
\toprule\noalign{}
\begin{minipage}[b]{\linewidth}\raggedright
Policy
\end{minipage} & \begin{minipage}[b]{\linewidth}\raggedleft
Scenario 1: EXP-009 shock
\end{minipage} & \begin{minipage}[b]{\linewidth}\raggedleft
Scenario 2: severe shock
\end{minipage} \\
\midrule\noalign{}
\endhead
\bottomrule\noalign{}
\endlastfoot
Oracle & 104,680 & 22,829 \\
\textbf{AlwaysRATED} (best constant) & \textbf{99,321} & \textbf{10,364} \\
LinTS & 71,429 (−39.0 \%, p \textless{} 0.0001, d = −4.3) & −14,296 (p \textless{} 0.0001, d = −4.6) \\
LinUCB & 71,373 (−39.2 \%, p \textless{} 0.0001, d = −5.3) & −14,443 (p \textless{} 0.0001, d = −5.7) \\
DiscountedLinUCB (\(\lambda = 0.999\)) & 70,995 (−39.9 \%) & −20,542 \\
DiscountedLinUCB (\(\lambda = 0.995\)) & 68,919 & --- \\
Static XGB & 52,715 & --- \\
\end{longtable}
}

\fg{0.85\textwidth}{Figure 17. EXP-015 total reward across the full ladder under the EXP-009 shock; AlwaysRATED (constant) dominates every learner. A severe-shock companion figure (fig\_drift\_rescue\_ladder\_severe, supplementary) shows the learners turning net-negative while the constant stays positive}{images/fig_drift_rescue_ladder.pdf}

\textbf{Verdict.} Under the realistic shock, \textbf{no adaptive policy beats the best constant} --- the gap is \textasciitilde39 \% with d ≈ −4 to −5. Under the severe shock the learners go \textbf{net-negative while the constant stays positive} --- exactly when the book is most marginal, the bandit's exploration cost is most punishing. \textbf{Forgetting does not help at any \(\lambda\).} The reason is structural: the shock reshuffles \emph{who} falls in each risk bucket but does not change the near-constant \emph{character} of the optimal policy (``rate most, decline the clearly-uninsurable''), so a fixed ``rate everyone'' policy is barely hurt and there is little for an adaptive learner to exploit. Non-stationarity --- the last and most natural place to seek the bandit's value --- does not rescue it.

\section{EXP-010: Cold-Start --- When Does the Bandit Catch Up?}\label{exp-010-cold-start-when-does-the-bandit-catch-up}

\fg{0.85\textwidth}{Figure 18. Cumulative reward at horizons T ∈ {200, 500, 1,000, 2,000} for LinUCB, LinTS, and a freshly-trained XGBoost baseline (mean of 10 seeds). FreshXGB leads at T ≤ 500; bandits cross over between T = 1,000 and T = 2,000}{images/fig_010_cold_start.pdf}

\subsection{Design}\label{design-3}

EXP-010 asks an operationally crucial question that the headline experiments do not: \emph{at what data scale does the bandit begin to beat a freshly-trained static model?} At each horizon \(T \in \{200, 500, 1{,}000, 2{,}000\}\), an XGBoost regressor is retrained from scratch on the first \(T\) applicants and then deployed under the same fixed rule engine as the headline Static XGB (the \textbf{FreshXGB} baseline). LinUCB and LinTS are run with their standard parameters (\(\alpha = 1.0\), \(v^2 = 1.0\)) for exactly \(T\) rounds, starting from \(A = I, b = 0\). Cumulative rewards at each horizon are compared across \textbf{10 seeds per horizon} owing to the cost of repeated XGBoost refits (\texttt{exp\_010\_cold\_start\_analysis.py:253}).

\subsection{Results}\label{results-4}

\captionof{table}{Table 21 --- EXP-010 cumulative reward by horizon (10 seeds, mean ± std)}\label{tab:5.10.1}

{\def\LTcaptype{none} % do not increment counter
\begin{longtable}[]{@{}
  >{\raggedleft\arraybackslash}p{(\linewidth - 8\tabcolsep) * \real{0.2105}}
  >{\raggedleft\arraybackslash}p{(\linewidth - 8\tabcolsep) * \real{0.2105}}
  >{\raggedleft\arraybackslash}p{(\linewidth - 8\tabcolsep) * \real{0.2105}}
  >{\raggedleft\arraybackslash}p{(\linewidth - 8\tabcolsep) * \real{0.2105}}
  >{\raggedright\arraybackslash}p{(\linewidth - 8\tabcolsep) * \real{0.1579}}@{}}
\toprule\noalign{}
\begin{minipage}[b]{\linewidth}\raggedleft
Horizon \(T\)
\end{minipage} & \begin{minipage}[b]{\linewidth}\raggedleft
LinUCB
\end{minipage} & \begin{minipage}[b]{\linewidth}\raggedleft
LinTS
\end{minipage} & \begin{minipage}[b]{\linewidth}\raggedleft
FreshXGB
\end{minipage} & \begin{minipage}[b]{\linewidth}\raggedright
Bandit winner?
\end{minipage} \\
\midrule\noalign{}
\endhead
\bottomrule\noalign{}
\endlastfoot
200 & 1,313 ± 971 & 1,206 ± 1,038 & \textbf{2,812 ± 1,115} & No --- FreshXGB ≈ 2.1× \\
500 & 5,593 ± 1,597 & 5,570 ± 1,051 & \textbf{7,048 ± 1,407} & No --- FreshXGB ≈ 1.26× \\
1,000 & 13,033 ± 2,174 & 13,496 ± 949 & 13,692 ± 1,819 & Tie (within seed variance) \\
2,000 & \textbf{31,227 ± 2,987} & \textbf{31,401 ± 2,187} & 29,267 ± 1,731 & \textbf{Yes --- bandits cross over} \\
\end{longtable}
}

At \(T = 200\) the freshly-trained XGBoost regressor more than doubles the bandits' cumulative reward. At \(T = 500\) it retains a one-quarter advantage. By \(T = 1{,}000\) the three algorithms are statistically indistinguishable. The \textbf{operational crossover} --- the smallest horizon at which the bandits beat a fresh static model --- falls between \(T = 1{,}000\) and \(T = 2{,}000\) on the Cambodia dataset.

\subsection{Interpretation and operational implications}\label{interpretation-and-operational-implications}

The cold-start result is the honest counterweight to the §5.1 finding --- but it must be read alongside §5.13. EXP-010 shows the bandit overtakes a \emph{freshly-trained XGBoost rule} between T = 1,000 and 2,000. It does \textbf{not} show the bandit overtaking the \textbf{constant}: §5.13 runs to 20,000 rounds and finds the bandit's cumulative reward \emph{never} crosses AlwaysRATED, and its converged per-round reward plateaus below it. So the crossover quantified here is the horizon at which online adaptation beats a \emph{frozen rule} --- an operational fact about warm-starting --- not a horizon at which the bandit becomes the best available policy. The headline experiment runs to \(N = 5{,}000\) rounds; the bandit has accumulated sufficient coefficient data to dominate the fixed-threshold baseline. At small \(T\) the situation reverses: a freshly trained XGBoost regressor, fitted offline on whatever historical data is available, exploits the dataset more efficiently than a bandit that begins each run from an uninformative prior.

The operational implication for a Cambodian insurer is direct: \textbf{a contextual bandit should not be deployed cold on the first 500--1,000 policy applications without warm-starting from an existing model}. Two warm-start strategies follow naturally from this finding and the live demo design (§4.10):

\begin{enumerate}
\def\labelenumi{\arabic{enumi}.}
\tightlist
\item
  \textbf{Pre-train then deploy.} Run a Static XGB or GLM baseline offline on historical data and use the predicted mortality multipliers to populate the bandit's \(b_a\) vectors with synthetic observations before the first live applicant arrives. The bandit then begins production already in the regime where it dominates.
\item
  \textbf{Shadow Mode warm-up.} Run the bandit alongside the existing rule engine for the first 1,000--2,000 production applications, updating its sufficient statistics from the static rule's outcomes (§4.10.1). Switch to bandit-binding decisions only after the crossover horizon is reached.
\end{enumerate}

Neither strategy is novel --- both are documented in the contextual-bandit deployment literature --- but EXP-010 quantifies for the first time the horizon at which the switch becomes net-positive on the Cambodia dataset specifically.

\subsection{Pass/Fail Assessment}\label{passfail-assessment-7}

\begin{enumerate}
\def\labelenumi{\arabic{enumi}.}
\tightlist
\item
  \textbf{LinUCB crossover exists (beats FreshXGB at T = 2,000).} \(\$31{,}227 > \$29{,}267\). PASSED.
\item
  \textbf{LinTS crossover exists (beats FreshXGB at T = 2,000).} \(\$31{,}401 > \$29{,}267\). PASSED.
\item
  \textbf{LinUCB cumulative reward \textgreater{} FreshXGB at T = 2,000.} Confirmed (criterion 1). PASSED.
\item
  \textbf{LinTS cumulative reward \textgreater{} FreshXGB at T = 2,000.} Confirmed (criterion 2). PASSED.
\end{enumerate}

EXP-010: PASS, with the operational deployment-boundary qualifier of §5.10.3.

\section{Reward-Model Sensitivity --- Is the Negative Result an Artifact of the Reward Design?}\label{reward-model-sensitivity}

The headline finding of this chapter --- that a one-line constant policy (AlwaysRATED) outperforms every deployable learner --- could in principle be an artifact of the particular reward function the simulator uses. This section tests that objection directly by re-scoring all eight policies under six reward models: the as-published \textbf{ORIGINAL} model; \textbf{A}, which symmetrises the adverse-selection response; \textbf{B}, which adds A to a symmetric claims-cost split; and \textbf{C}, a constructed fair-pricing regime swept at three premium loadings (\(\ell = 0.45, 0.60, 0.75\)). Each policy is evaluated on the same 2,000-applicant population over 20 seeds under common random numbers, so the only thing that changes across columns is the reward design. Reward, not regret, is the comparison metric: for the constant policies the realized-regret sign is distorted by claims noise, so the common-random-number reward --- whose Oracle, 126,351, agrees with the EXP-007 harness to within rounding --- is the like-for-like quantity.

\subsection{Results}\label{reward-sensitivity-results}

\captionof{table}{Table 22 --- Mean cumulative reward by policy and reward model (20 seeds, common random numbers). A = symmetric adverse-selection; B = A + symmetric cost; C = fair-priced at loading \(\ell\).}\label{tab:5.11.1}

{\def\LTcaptype{none}
\footnotesize
\begin{longtable}{@{}lrrrrrr@{}}
\toprule
Policy & ORIGINAL & A & B & C $\ell$=.45 & C $\ell$=.60 & C $\ell$=.75 \\
\midrule
\endhead
\bottomrule
\endlastfoot
Oracle & 126,351 & 108,690 & 98,537 & 85,749 & 109,099 & 127,651 \\
LogisticOracle\textsuperscript{\ddag} & 123,977 & 107,186 & 97,182 & 83,479 & 102,176 & 116,248 \\
\textbf{AlwaysRATED} & \textbf{122,287} & \textbf{99,524} & \textbf{80,615} & \textbf{82,229} & \textbf{103,391} & \textbf{116,922} \\
LinTS & 93,723 & 79,122 & 66,426 & 62,748 & 76,432 & 89,085 \\
LinUCB & 91,864 & 75,866 & 65,589 & 61,223 & 76,483 & 88,574 \\
StaticXGB\textsuperscript{\dag} & 72,206 & 56,297 & 41,757 & 45,844 & 56,234 & 63,153 \\
AlwaysSTANDARD & 34,684 & 34,684 & 17,955 & 50,819 & 50,819 & 50,819 \\
Random & 1,980 & -6,543 & -15,447 & -10,690 & -1,815 & 4,149 \\
\end{longtable}
}

\textsuperscript{\dag}\,StaticXGB's actions are model-independent (thresholds tuned for ORIGINAL); only its reward changes per model. \textsuperscript{\ddag}\,LogisticOracle is a frozen classifier trained in-sample on the oracle-optimal labels --- an optimistic ceiling, not a deployable baseline.

The policy ordering is \textbf{invariant across all six reward models}: Oracle \(>\) AlwaysRATED \(>\) LinTS \(\approx\) LinUCB \(>\) StaticXGB \(>\) AlwaysSTANDARD \(>\) Random. AlwaysRATED is the best \emph{deployable} policy in every model, and the two bandits sit a full tier below it --- holding only 74--82\% of the best constant's reward regardless of the reward design. The 95\% bootstrap confidence intervals of AlwaysRATED and the best deployable learner are disjoint in every model (ORIGINAL: AlwaysRATED \([119{,}490, 125{,}084]\) versus LinTS \([91{,}162, 96{,}447]\)).

\captionof{table}{Table 23 --- Each learner versus AlwaysRATED (paired Wilcoxon over 20 seeds; Cohen's \(d\)).}\label{tab:5.11.2}

{\def\LTcaptype{none}
\begin{longtable}{@{}lrrr@{}}
\toprule
Reward model & \(p\) (both learners) & \(d\) (LinUCB) & \(d\) (LinTS) \\
\midrule
\endhead
\bottomrule
\endlastfoot
ORIGINAL & \(< 0.0001\) & -5.14 & -4.51 \\
A (sym-adverse) & \(< 0.0001\) & -4.79 & -4.59 \\
B (A + sym-cost) & \(< 0.0001\) & -3.47 & -4.23 \\
C ($\ell$=.45) & \(< 0.0001\) & -6.94 & -6.58 \\
C ($\ell$=.60) & \(< 0.0001\) & -5.87 & -7.46 \\
C ($\ell$=.75) & \(< 0.0001\) & -5.23 & -5.83 \\
\end{longtable}
}

Every learner is significantly worse than AlwaysRATED in every model, with large-to-enormous negative effect sizes (\(|d| = 3.5\)--\(7.5\)). This is not a marginal loss that a reward re-specification might erase.

\subsection{Interpretation}\label{reward-sensitivity-interpretation}

Two things move with the reward design and one decisive thing does not. What moves: the \emph{absolute} reward levels (every correction lowers them) and AlwaysRATED's \emph{share of the oracle}, which falls from 96.8\% under ORIGINAL to 81.8\% under model B --- confirming that part of the constant's original dominance was an adverse-selection and cost artifact, and that the earlier ``97.6\% of oracle'' figure overstated how near-optimal blanket rating is. What does \emph{not} move: the policy ranking, and the bandit's fraction of the best constant (74--82\% throughout). The two quantities a reviewer cares about most are exactly the ones that are insensitive to reward design.

Crucially, model C is built to \emph{favour} segmentation --- it removes the pricing free-lunch and lets the oracle diversify (the optimal STANDARD share rises with the loading) --- yet AlwaysRATED still wins at all three loadings. And because a fully-supervised linear model (LogisticOracle) tracks the oracle in every model, the optimal policy \emph{is} linearly representable in the bandit's own features: the bandit's shortfall is a partial-feedback and exploration cost, not an absence of learnable structure. The dominance of the constant is therefore a robust property of the learning problem, not an artifact of any single reward asymmetry.

\section{EXP-016: Demand-Elasticity Sweep --- Does Any Demand Regime Rescue the Bandit?}\label{exp-016-demand-elasticity-sweep}

EXP-012's adverse-selection and elasticity checks (§5.8) varied demand only as a robustness probe of the LinUCB-versus-Static gap. This section runs the \emph{dedicated} experiment that the root-cause analysis demanded: a seven-point sweep of demand elasticity comparing each learner against the \emph{best constant} --- the one reward component never previously swept against the full ladder. Demand elasticity is varied through the acceptance slope, with the acceptance intercept re-anchored so that mean acceptance stays at \(\approx 0.5\) at the population-mean price-to-income ratio for every slope, keeping the demand curve non-degenerate at all points. This is a sweep that maps where contextual value emerges, not a search for a slope that flatters the bandit. For reference, real low- and middle-income-country microinsurance demand is typically inelastic (\(\approx -0.2\) to \(-0.5\)); the simulator's own demand sits in the mildly-elastic range.

\subsection{Results}\label{elasticity-results}

\captionof{table}{Table 24 --- Demand-elasticity sweep: mean cumulative reward (20 seeds). No learner beats the best constant at any elasticity. Const/Oracle = best constant as a percentage of the oracle.}\label{tab:5.12.1}

{\def\LTcaptype{none}
\footnotesize
\begin{longtable}{@{}lrrrrrr@{}}
\toprule
Elasticity & Best const. & LinUCB & LinTS & Log.\,Oracle & Oracle & Const/Oracle \\
\midrule
\endhead
\bottomrule
\endlastfoot
\(-0.24\) & 365,385 & 281,329 & 284,289 & 365,385 & 365,385 & 100.0\% \\
\(-0.45\) & 337,763 & 266,141 & 259,913 & 337,786 & 338,215 & 99.9\% \\
\(-0.78\) & 309,207 & 244,762 & 247,180 & 309,095 & 310,120 & 99.7\% \\
\(-1.08\) & 299,935 & 248,440 & 249,497 & 302,535 & 304,207 & 98.6\% \\
\(-1.21\) & 302,980 & 255,610 & 257,789 & 308,654 & 311,419 & 97.3\% \\
\(-1.29\) & 307,886 & 264,503 & 263,361 & 317,927 & 323,556 & 95.2\% \\
\(-1.32\) & 311,532 & 265,938 & 273,091 & 326,599 & 336,514 & 92.6\% \\
\end{longtable}
}

At no elasticity does any learner beat the best constant: the bandits hold 77--88\% of the best constant's reward across the entire range. The sweep nonetheless contains a sharp, honest nuance rather than a flat null.

\subsection{Interpretation}\label{elasticity-interpretation}

At the realistic, inelastic end (\(-0.24\)) the best constant \emph{is} the oracle: AlwaysRATED, LogisticOracle, and the oracle all return the identical 365,385, the optimal action is RATED for every income group, and there is no contextual structure to exploit --- so the bandit is most superfluous exactly where the application actually lives. As demand becomes more elastic, genuine contextual structure emerges: the oracle begins to split on income (at the most-elastic cell, low-income applicants become optimally STANDARD or DECLINE while high-income applicants stay RATED), and the gap between the best constant and the oracle widens from 0\% to 7.4\%. A fully-supervised linear model captures that structure --- LogisticOracle tracks the oracle at 97--100\% across the whole range, including the elastic end --- confirming the structure is present and linearly representable in the bandit's own features. Yet the contextual bandit never captures it: even where the structure is richest, LinTS (273,091) sits below the best constant (311,532), not above. The limiting factor is learning the income-split policy from scalar, low-acceptance bandit feedback while paying exploration cost --- not the absence of structure.

Combined with the reward-model sensitivity of §5.11, the drift rescue of §5.9.5, and the HITL re-benchmark of §5.4, the elasticity sweep closes the last axis: there is no tested regime --- reward model, non-stationarity, feedback structure, or demand elasticity --- in which the contextual bandit beats a trivial constant policy. Notably, a \emph{more realistic, more inelastic} demand curve makes the constant \emph{more} dominant, not less, so demand modelling cannot rescue the contribution.

\section{Horizon Check --- Is the Defeat Structural or Cold-Start Cost?}\label{horizon-check}

The headline experiments run for 5,000 rounds (\(\approx 2.5\) passes through the 2,000-applicant population), which charges the bandit its full cold-start learning cost amortised over a comparatively short horizon. A natural objection to the negative result is therefore that the bandit simply needs more data: given enough rounds, its converged policy might overtake the constant. This section settles that objection by running a long horizon --- 20,000 rounds (ten passes) under the ORIGINAL reward, five seeds --- and asking two discriminating questions: does the bandit's \emph{steady-state} per-round reward (the last 1,000 rounds, after convergence) reach the constant's, and does its \emph{cumulative} reward ever cross AlwaysRATED?

\captionof{table}{Table 25 --- Per-round reward by window over a 20,000-round horizon (5 seeds, ORIGINAL reward).}\label{tab:5.13.1}

{\def\LTcaptype{none}
\begin{longtable}{@{}lrrrr@{}}
\toprule
Policy & Rounds 0--1k & 4--5k & 9--10k & 19--20k \\
\midrule
\endhead
\bottomrule
\endlastfoot
Oracle & 26.61 & 25.27 & 24.08 & 23.98 \\
AlwaysRATED & 25.26 & 25.24 & 23.88 & 22.75 \\
LinUCB & 13.79 & 21.05 & 20.38 & 20.47 \\
LinTS & 13.49 & 21.26 & 21.64 & 20.95 \\
\end{longtable}
}

The answer to both questions is no. After convergence the bandits' steady-state per-round reward plateaus at 20.5 (LinUCB) and 21.0 (LinTS), both \emph{below} the constant's 22.8 and well below the oracle's 24.0. And cumulative reward never crosses AlwaysRATED within 20,000 rounds: the final cumulative shortfall is \$70,324 for LinUCB and \$61,546 for LinTS, and it is still widening, not closing, at the end of the run. The bandit's defeat is therefore \textbf{structural}, not a finite-horizon cold-start artifact --- more data does not rescue it.

This resolves an apparent tension with EXP-013 (§5.6). EXP-013 validated that LinUCB's regret grows at the \(\tilde O(d\sqrt{T})\) rate --- but that is the rate at which the bandit converges to its \emph{own} optimum. The horizon check shows that this optimum is itself below the constant, so the regret tail measured against the \emph{best simple policy} is ultimately linear, not \(\sqrt{T}\): a \(\sqrt{T}\)-consistent learner converging to a suboptimal target never overtakes the target it should have matched.

